% Options for packages loaded elsewhere
\PassOptionsToPackage{unicode}{hyperref}
\PassOptionsToPackage{hyphens}{url}
\documentclass[
]{article}
\usepackage{xcolor}
\usepackage{amsmath,amssymb}
\setcounter{secnumdepth}{-\maxdimen} % remove section numbering
\usepackage{iftex}
\ifPDFTeX
  \usepackage[T1]{fontenc}
  \usepackage[utf8]{inputenc}
  \usepackage{textcomp} % provide euro and other symbols
\else % if luatex or xetex
  \usepackage{unicode-math} % this also loads fontspec
  \defaultfontfeatures{Scale=MatchLowercase}
  \defaultfontfeatures[\rmfamily]{Ligatures=TeX,Scale=1}
\fi
\usepackage{lmodern}
\ifPDFTeX\else
  % xetex/luatex font selection
\fi
% Use upquote if available, for straight quotes in verbatim environments
\IfFileExists{upquote.sty}{\usepackage{upquote}}{}
\IfFileExists{microtype.sty}{% use microtype if available
  \usepackage[]{microtype}
  \UseMicrotypeSet[protrusion]{basicmath} % disable protrusion for tt fonts
}{}
\makeatletter
\@ifundefined{KOMAClassName}{% if non-KOMA class
  \IfFileExists{parskip.sty}{%
    \usepackage{parskip}
  }{% else
    \setlength{\parindent}{0pt}
    \setlength{\parskip}{6pt plus 2pt minus 1pt}}
}{% if KOMA class
  \KOMAoptions{parskip=half}}
\makeatother
\usepackage{longtable,booktabs,array}
\newcounter{none} % for unnumbered tables
\usepackage{calc} % for calculating minipage widths
% Correct order of tables after \paragraph or \subparagraph
\usepackage{etoolbox}
\makeatletter
\patchcmd\longtable{\par}{\if@noskipsec\mbox{}\fi\par}{}{}
\makeatother
% Allow footnotes in longtable head/foot
\IfFileExists{footnotehyper.sty}{\usepackage{footnotehyper}}{\usepackage{footnote}}
\makesavenoteenv{longtable}
\setlength{\emergencystretch}{3em} % prevent overfull lines
\providecommand{\tightlist}{%
  \setlength{\itemsep}{0pt}\setlength{\parskip}{0pt}}
\usepackage{bookmark}
\IfFileExists{xurl.sty}{\usepackage{xurl}}{} % add URL line breaks if available
\urlstyle{same}
\hypersetup{
  hidelinks,
  pdfcreator={LaTeX via pandoc}}

\author{}
\date{}

\begin{document}

\section{Home@ix '' Full Working Book
(Draft)}\label{homeix-full-working-book-draft}

\textbf{Print bundle:} print\_20260302\_230225 Â~\textbar Â~
\textbf{Built:} 2026-03-02 23:03

\subsection{Known issues
(pre-publication)}\label{known-issues-pre-publication}

\begin{itemize}
\tightlist
\item
  \textbf{Animation rendering:} some animated assets may not display in
  all viewers. Use exported GIF/PNG fallbacks in \texttt{assets/} where
  provided.
\item
  \textbf{Math rendering:} equations are standardised to
  \texttt{\$\$...\$\$}. If any equations appear unrendered, open the
  HTML in a modern browser (MathJax-enabled).
\end{itemize}

\begin{center}\rule{0.5\linewidth}{0.5pt}\end{center}

\protect\phantomsection\label{my-new-section-title}
\protect\phantomsection\label{chapter-1}
\subsection{Chapter 1 '' Introduction: The Regime Behind the
Crisis}\label{chapter-1-introduction-the-regime-behind-the-crisis}

The UK housing affordability crisis is routinely described as a shortage
story. Not enough homes. Not enough land released. Not enough planning
permissions granted. That framing is not wrong '' but it is radically
incomplete. And its incompleteness has policy consequences: it directs
attention toward symptoms while the structural engine of exclusion runs
largely unexamined.

This book argues that the dominant driver of accelerating homeownership
exclusion in the UK is not a physical supply gap alone. It is a
\textbf{financial regime failure} '' one shaped by the concentration of
bank credit in residential mortgage lending, the macroprudential
architecture of post-crisis banking regulation (Basel III), and the
resulting feedback loops that make housing markets clear through
\emph{rationing and exclusion} rather than through orderly repricing and
widening access.

\textbf{Core thesis:} When bank balance sheets are structurally
dominated by residential mortgage assets, and when macroprudential
regulation tightens capital and underwriting standards in response to
systemic risk, the transmission mechanism runs directly through housing
access. Credit tightens. Turnover collapses. The transacting population
narrows. Mortgage-dependent households '' disproportionately younger,
lower-wealth, and renting '' are screened out. The observed price index
becomes a \emph{selected statistic} generated by a thinner, wealthier
pool of buyers. Affordability deteriorates not because prices spike
dramatically, but because the market stops being accessible at all.

\begin{center}\rule{0.5\linewidth}{0.5pt}\end{center}

\subsubsection{Part One: Financialisation as a Structural
Condition}\label{part-one-financialisation-as-a-structural-condition}

The UK banking sector allocates an unusually high share of
private-sector lending to residential mortgages '' a structural feature
that links housing-market cycles directly to bank balance-sheet health,
deposit creation, and broad money dynamics. In an endogenous money
system, mortgage lending \emph{creates} deposits: when mortgage
origination weakens, broad money growth can undershoot nominal activity
trends, tightening conditions economy-wide. When it surges, it inflates
collateral values, reinforcing the next cycle of lending against the
same asset base.

This is not a market functioning normally under price signals. It is a
\textbf{financialised circuit} in which housing serves simultaneously as
necessity, collateral, and speculative asset '' and in which the
interests of existing asset-holders are structurally privileged over the
access needs of would-be entrants. The consequence is not merely high
prices. It is a market that, under stress, stops clearing for a large
share of the population entirely.

\begin{center}\rule{0.5\linewidth}{0.5pt}\end{center}

\subsubsection{Part Two: Basel III and the Macroprudential
Transmission}\label{part-two-basel-iii-and-the-macroprudential-transmission}

The post-Global Financial Crisis regulatory settlement '' crystallised
in the Basel III framework and its successive implementations '' was
designed to make banks safer. In important respects, it has. Capital
buffers are larger. Loss-absorbency requirements are more stringent.
Underwriting standards have tightened.

But the framework is, by design and by the admission of its critics,
\textbf{almost exclusively microprudential}: focused on the solvency of
individual institutions rather than the systemic allocation of credit
across the economy. It does not directly constrain the
\emph{composition} of bank lending portfolios '' the structural
overweight in residential mortgages that makes the financial system so
sensitive to housing-market shocks. And when tighter capital
requirements and more conservative underwriting are applied to a system
already dominated by mortgage credit, the transmission is asymmetric: it
falls hardest on the marginal borrower.

The marginal borrower in the UK housing market is, overwhelmingly, the
first-time buyer: younger, lower-equity, more dependent on high
loan-to-value products, and most exposed to the deposit and
income-multiple constraints that tighten when underwriting standards
rise. Basel III's capital surcharges, countercyclical buffers, and
stress-testing regimes are not designed to exclude this cohort. But in a
financialised housing system, that is their operational effect.
\textbf{Regulatory prudence and structural exclusion become, in
practice, the same mechanism.}

\begin{center}\rule{0.5\linewidth}{0.5pt}\end{center}

\subsubsection{Part Three: FAIR as the Empirical and Testable Evidence
Base}\label{part-three-fair-as-the-empirical-and-testable-evidence-base}

Diagnosis without measurement is assertion. The third and central
contribution of this book is the introduction of \textbf{Home@ix FAIR}
'' a transparent, reproducible, quarterly indicator designed to detect
when the UK housing system is shifting into a regime of accelerated
affordability deterioration, understood not as a price ratio but as an
\textbf{access-and-allocation outcome}.

FAIR is built from two observable precursors to exclusion stress that
are directly legible in aggregate data:

\begin{itemize}
\tightlist
\item
  \textbf{Credit``price decoupling:} house-price growth outrunning
  mortgage-stock growth '' the signal that prices are detaching from the
  credit base that finances them, compressing access for
  mortgage-dependent buyers.
\item
  \textbf{Thin-market clearing:} turnover declining relative to stock ''
  the signal that the market is rationing participation rather than
  clearing through price adjustment.
\end{itemize}

Together, these two mechanisms encode the financialisation and
macroprudential transmission story in a form that is \emph{quarterly,
auditable, and backtestable}. When credit tightens under Basel III-era
underwriting standards, the wedge between price growth and
mortgage-stock growth widens. When the eligible buyer pool narrows,
turnover falls. FAIR rises. The regime signal fires '' typically 2``3
years before the exclusion effect is visible in conventional
affordability statistics.

\textbf{What FAIR adds to the literature:} existing affordability
measures are price-centric and contemporaneous. FAIR is regime-centric
and forward-looking. It does not ask ``are prices high?'' It asks: ``is
the system moving into a state where access is about to deteriorate, and
is that deterioration being driven by credit-channel dynamics rather
than by fundamental demand?'' That is a different '' and, this book
argues, more policy-relevant '' question.

\begin{center}\rule{0.5\linewidth}{0.5pt}\end{center}

\subsubsection{Structure of This Book}\label{structure-of-this-book}

The book is organised as a working research bundle. Each section is
self-contained but the argument builds sequentially:

\begin{itemize}
\tightlist
\item
  \textbf{Chapter 1} (this chapter) frames the financial regime failure,
  the Basel III transmission mechanism, and FAIR as the empirical spine.
\item
  \textbf{Chapter 2} develops the four interlocking dimensions of crisis
  '' credit impairment, mortgage concentration, M4 money supply
  dependency, and affordability as allocation failure '' and shows how
  they amplify each other through coupled feedback loops.
\item
  \textbf{Chapter 3} examines why standard policy responses cannot break
  the feedback loops, and sets out the structural reform requirements
  implied by the regime diagnosis.
\item
  \textbf{Part I and the FAIR Papers} that follow provide the
  foundational market-function exhibits, the formal FAIR derivation, and
  the full reproducible technical specification.
\end{itemize}

The figures throughout are generated directly from the Home@ix data
pipeline. Every number is traceable to a canonical input. Every chart is
regenerable from the accompanying scripts. This is not only a research
argument '' it is a \textbf{reproducible evidence base}, designed to be
stress-tested, extended, and challenged.

\textbf{An invitation:} the FAIR indicator is intentionally open. Its
weights are fixed by design, not estimated '' which means they can be
varied and the sensitivity tested. Its baseline windows are explicit ''
which means they can be contested. Its crash-start definitions are
rule-based '' which means they can be replicated by anyone with the
data. The goal is not to close the debate about UK housing
financialisation and regulatory transmission. It is to give that debate
a reproducible empirical floor.

\begin{center}\rule{0.5\linewidth}{0.5pt}\end{center}

\protect\phantomsection\label{chapter-2}
\subsection{Chapter 2 '' The Four Dimensions: An Integrated
Analysis}\label{chapter-2-the-four-dimensions-an-integrated-analysis}

The UK housing system is not suffering from a single, tractable problem.
It is suffering from four interlocking structural failures that
reinforce each other through feedback loops. Analysed separately '' as
they usually are '' each dimension looks manageable. Understood as a
coupled system, they describe a regime in which crisis is not a tail
risk but a probable trajectory. This chapter develops each dimension in
turn before showing how they combine.

\textbf{The four dimensions:}

\begin{enumerate}
\def\labelenumi{\arabic{enumi}.}
\tightlist
\item
  Credit impairment risk in commercial banks under mortgage
  concentration
\item
  Mortgage lending over-reliance and the collateral feedback loop
\item
  Structural dependency of broad money (M4) on housing credit creation
\item
  Affordability deterioration as an access-and-allocation regime outcome
\end{enumerate}

\begin{center}\rule{0.5\linewidth}{0.5pt}\end{center}

\subsubsection{2.1 Dimension One: Credit Impairment '' Surface Stability
Masking Structural
Fragility}\label{dimension-one-credit-impairment-surface-stability-masking-structural-fragility}

The UK mortgage market presents a paradoxical picture. Headline arrears
stood at approximately 0.9\% of outstanding mortgages in Q1 2025 ''
historically low, and apparently healthy relative to pre-GFC peaks of
2``3\%. Yet this surface benignancy obscures three significant
vulnerabilities.

First, the low arrears rate reflects \textbf{compositional bias}: as
mortgage rates declined from their 2022``23 peak, better-quality
borrowers re-entered the market while distressed cohorts had already
exited or been screened out through stricter underwriting. The observed
pool is disproportionately high-quality, understating true vulnerability
in the broader stock. Second, the published arrears threshold (2.5\% of
balance) captures only serious delinquency, missing earlier-stage stress
signals that predict acceleration. Third, serious delinquencies were
already rising: the rate reached 1.0\% in September 2025, up from 0.9\%
year-on-year '' an early stress signal emerging beneath the headline.

\paragraph{Buy-to-Let Arrears: The Leading
Indicator}\label{buy-to-let-arrears-the-leading-indicator}

Buy-to-let mortgage arrears reached 2.5\% in Q1 2025 '' approximately
2.8 times the owner-occupied rate. This divergence reflects the
structural exposure of the BTL cohort: landlord borrowers have no
fixed-income anchor against rising rates, and their cashflows are
directly exposed to tenant financial deterioration. Approximately one in
five UK households now rents privately; as renter financial
vulnerability rises, tenant payment defaults increase, which in turn
pressures landlord mortgage payments. BTL arrears have historically
preceded owner-occupied arrears by 6``12 months, making the current BTL
rate a meaningful forward signal for the broader system.

\paragraph{The Debt-Servicing Ratio: A Linear, Not Threshold,
Risk}\label{the-debt-servicing-ratio-a-linear-not-threshold-risk}

Bank Underground research (2024) established that the relationship
between mortgage debt-servicing ratios (DSR) and arrears probability is
smooth and linear '' not a discrete cliff above some critical threshold.
This has profound implications: \textbf{deterioration across the entire
DSR distribution matters for financial stability}, not merely tail risk.
The current distribution places approximately 40\% of mortgagors below
15\% DSR (low stress), approximately 50\% between 15``30\% (manageable
but vulnerable), approximately 7``10\% between 30``40\% (elevated
stress), and approximately 2``5\% above 40\% (high stress). A modest
rightward shift of 2``3 percentage points across the distribution ''
conceivable under employment stress '' would move a significant share of
the mortgage population into elevated-stress territory.

The unemployment transmission is direct: a 1 percentage point rise in
unemployment historically correlates with a 30``50\% increase in
mortgage arrears. At current unemployment of approximately 3.8\%, a rise
to 5\% would push arrears from 0.9\% toward 1.5``2.0\%. The capital
impact is material: system-wide mortgage losses of £10 billion would
reduce bank lending capacity by approximately £90 billion at the
regulatory 11\% Tier 1 minimum.

\paragraph{High-LTV Lending Expansion: A Recent Negative
Trend}\label{high-ltv-lending-expansion-a-recent-negative-trend}

A significant structural shift has emerged since 2024. First-time buyer
average LTV reached 77.3\% '' a 15-year high. The share of gross
mortgage advances with LTV exceeding 90\% rose from 6.7\% (Q3 2024) to
7.1\% (Q2 2025). The supply of 95\% LTV mortgages reached approximately
18-year highs in February 2026. This expansion reflects both easing
regulatory constraints and competitive pressure among lenders.

The financial stability implications are severe. Loss-given-default
(LGD) rises sharply with LTV: below 60\% LTV, LGD is typically 5``10\%;
at 80``90\% LTV, it rises to 20``30\%; above 90\%, it can exceed
40``50\%. A house price decline of 15\% '' historically plausible in
recession '' would push mortgages with LTV above 85\% into negative
equity. Underwater mortgages show 5``10 times higher default rates than
mortgages with equity cushion, creating potential for cascade effects if
price declines are sufficient to generate substantial negative equity
populations.

\paragraph{Basel 3.1 and the Regulatory
Paradox}\label{basel-3.1-and-the-regulatory-paradox}

The PRA's implementation of Basel 3.1 (effective January 2026, full
phase-in by 2030) recalibrates how banks measure risk and allocate
capital. Under the standardised approach, owner-occupied mortgage risk
weights range from 20``70\% depending on LTV; buy-to-let and
cash-flow-dependent properties attract 45``70\%+. Higher capital
requirements per pound of mortgage lending create return-on-equity
pressure, forcing banks to tighten underwriting standards, reduce
average loan sizes, or reduce lending volumes.

The result is a policy paradox: macroprudential regulation intended to
prevent defaults by improving bank resilience actually \textbf{causes}
defaults by reducing money creation and destroying demand in the real
economy. Credit rationing begins before defaults occur. The mechanism is
not theoretical '' it is the operational transmission channel between
regulatory tightening and housing exclusion.

\begin{center}\rule{0.5\linewidth}{0.5pt}\end{center}

\subsubsection{2.2 Dimension Two: Mortgage Lending Concentration and the
Collateral Feedback
Loop}\label{dimension-two-mortgage-lending-concentration-and-the-collateral-feedback-loop}

Residential mortgages comprise approximately 60\% of all UK bank lending
to the private sector '' approximately £1.73 trillion of the £2.87
trillion total MFI sterling lending stock. This concentration is not a
market equilibrium; it is a post-GFC anomaly driven by regulatory
incentives. Pre-2008, housing represented 50``55\% of bank credit. The
surge to 60\%+ reflects the capital efficiency advantage that Basel
risk-weighting confers on mortgage lending relative to business lending:
residential mortgages attract 20``70\% risk weights versus 65``135\% for
corporate exposures. For a capital-constrained bank, the optimal
allocation is mortgage-heavy.

Comparative international data confirms the UK's exceptional position:
Canada exhibits 55``60\% mortgage concentration; the Eurozone averages
45``50\%; the United States, 40``45\%. The UK's 60\%+ concentration
means housing-market shocks transmit directly and immediately to bank
balance sheets '' a structural vulnerability that no amount of
individual bank resilience can fully offset.

Business lending has been the casualty. Corporate lending fell from
approximately 30``35\% of total bank credit pre-GFC to approximately
20``25\% currently. SME lending has been particularly weak: the
proportion of SMEs actively seeking credit fell from approximately 25\%
to approximately 15\% post-crisis, with an increasing share turning to
non-bank sources. The result is a two-speed credit system in which
residential mortgage flows readily while productive business investment
faces structural constraint.

\paragraph{The Absorption Rate Constraint: Supply Cannot
Compensate}\label{the-absorption-rate-constraint-supply-cannot-compensate}

The absorption rate '' the rate at which newly built homes can be sold
without materially disturbing market price '' represents the binding
near-term constraint on UK housing supply. Major housebuilders maintain
sales rates within a relatively narrow band: approximately 0.45``0.8
homes per sales outlet per week across the cycle, currently 0.5``0.6 on
a forward-looking basis. With approximately 3,500 active sales outlets
across England, this implies annual new home sales of approximately
90,000``110,000 '' against a documented need of 340,000 homes annually.
The deficit is approximately 63\%.

This constraint is not a physical limit. It reflects \textbf{developer
incentive structures}: margin-per-unit optimisation takes priority over
volume. Developers pace sales to protect prices; releasing too many
homes too quickly would depress the very asset values that underpin
their balance sheets and future land options. The absorption rate is
therefore invariant to standard supply-side policy '' a point developed
fully in Chapter 3.

Savills' research documents that the number of sites gaining planning
consent fell 31\% over the five years to 2023, even where total
consented plot volumes held relatively stable. Consent approvals are
increasingly concentrated on fewer, larger sites, constraining outlet
growth and limiting opportunities for smaller builders and new entrants.
The pipeline of outlets is narrowing even as planning reform discussions
continue.

\begin{center}\rule{0.5\linewidth}{0.5pt}\end{center}

\subsubsection{2.3 Dimension Three: M4 Money Supply '' Structural
Dependency on Housing
Credit}\label{dimension-three-m4-money-supply-structural-dependency-on-housing-credit}

In the modern UK banking system, money is created through the act of
lending, not through the multiplication of central bank reserves. When a
commercial bank approves a mortgage, it simultaneously creates a loan
asset and a deposit liability '' new purchasing power. The Bank of
England's 2014 paper \emph{Money Creation in the Modern Economy}
formalised this: ``Rather than banks lending out deposits that are
placed with them, the act of lending creates deposits '' the reverse of
the sequence typically described in textbooks.''

Because residential mortgages comprise 60\% of all bank credit,
\textbf{the M4 money supply has become structurally dependent on housing
credit originations}. When mortgage volumes are strong, M4 grows. When
mortgage volumes weaken, M4 growth stalls. Unlike diversified credit
systems where weakness in housing can be offset by strength in business
lending, the UK system has no such offset: business lending is weak and
non-cyclical; consumer credit is volatile and small; government lending
is policy-dependent. Housing cycles dominate M4 dynamics.

\paragraph{Current Monetary
Insufficiency}\label{current-monetary-insufficiency}

As of October 2025, M4 growth ran at approximately 2.5``3.5\% annually
'' notably below the UK's nominal GDP growth trend of 4``4.5\%. Monthly
net lending to the private sector stood at £13.0 billion, down from
£19.6 billion in September '' a 32\% month-on-month deceleration. The
composition reveals the constraint: household net borrowing totalled
£5.4 billion (£4.3 billion mortgages, £1.1 billion consumer credit),
while private non-financial corporations actively repaid a net £4.8
billion. Money creation is concentrated on household housing collateral;
corporate balance sheets are deleveraging. If household mortgage demand
weakens, there is no offsetting corporate demand. The system must slow.

Over three years, a shortfall of 1``2 percentage points below trend
represents approximately 4\% of GDP '' approximately £120 billion of
missing money supply. In a high-debt economy, nominal income growth
shortfalls translate into rising real debt burdens: fixed debt on
falling nominal income triggers defaults even without unemployment
increases.

\paragraph{M4 Velocity Collapse: The Hidden
Signal}\label{m4-velocity-collapse-the-hidden-signal}

M4 velocity '' nominal GDP divided by M4 stock '' has declined from
1.1``1.2 pre-GFC to approximately 0.90``0.95 currently: a 15``20\% fall.
This indicates households are accumulating precautionary savings,
perceiving future scarcity. October 2025 data confirms this:
interest-bearing sight deposits rose £5.5 billion and ISAs rose £4.2
billion in a single month. In the accounting identity, nominal GDP = M4
Ã--- Velocity. If M4 grows 2.5\% and velocity declines 1\% annually,
nominal GDP growth is 1.5\% '' well below the 4\% trend required to
service existing debt loads without deterioration.

\begin{center}\rule{0.5\linewidth}{0.5pt}\end{center}

\subsubsection{2.4 Dimension Four: Affordability as Allocation
Failure}\label{dimension-four-affordability-as-allocation-failure}

The UK housing affordability crisis presents a fundamental paradox:
official statistics describe a severe housing shortage, yet substantial
vacant and underutilised stock exists. Government targets 300,000 new
homes annually; research estimates need at 340,000; current supply is
approximately 233,000. Yet simultaneously, approximately 700,000+
derelict and long-term empty properties exist, alongside 500,000+ second
homes and millions of underoccupied dwellings. The system contains
1.5``2.0 million potentially available homes.

\textbf{This paradox resolves when understood as an allocation crisis,
not a shortage crisis.} Prices at 8``10 times earnings signal severe
scarcity. Yet turnover has collapsed from 9\% of stock annually pre-GFC
to 3.7\% currently '' a 59\% decline. In a functioning market clearing
via price, turnover should be inversely correlated with
price-to-earnings ratios: higher prices reduce quantity demanded,
reducing turnover. The UK shows both extremely high prices and extremely
low turnover simultaneously '' the signature of a market clearing by
\textbf{rationing}, not by price adjustment.

\paragraph{Compositional Clearing: The Cash-Share
Signal}\label{compositional-clearing-the-cash-share-signal}

The rising cash share of transactions is the clearest empirical signal
of compositional clearing. HM Land Registry data (available from March
2012) shows cash buyers have represented a stable 30``35\% of all
transactions since that date '' approximately 35\% as of February 2026.
Pre-GFC, cash buyers represented sub-10\% of transactions. This
structural shift indicates that as prices rose from 5Ã--- to 8``10Ã---
earnings, mortgaged buyers '' particularly first-time buyers and younger
cohorts '' were progressively rationed out, and cash-capable buyers
filled the vacuum. The market does not clear by price adjustment
reducing quantity demanded uniformly; it clears by excluding certain
buyer cohorts entirely.

\paragraph{The Verified Demand Trap}\label{the-verified-demand-trap}

Banks underwrite mortgages using income multiples of 4.5``5.5Ã--- and
stress-tested affordability. At 4.5Ã--- income and a 6\% effective
stress rate (4\% mortgage rate plus 2\% buffer), a £50,000-income
household can borrow up to £225,000 '' implying a debt-servicing ratio
of approximately 33\% of gross income at maximum tolerance. Yet
effective household demand '' the amount a household can sustainably
borrow while maintaining living standards and savings '' is typically
30\% lower. The same household budgeting at 20\% of after-tax income for
housing can sustain a mortgage of approximately £110,000.

The gap between verified demand (£225,000) and effective demand
(£110,000) is the engine of long-run price inflation. Prices are bid up
to the verified ceiling; effective-demand buyers are rationed out; the
marginal transacting buyer is the one leveraging to the maximum approved
amount. Over decades, this mechanism drove the price-to-earnings ratio
from approximately 3Ã--- in 1980 to approximately 8``10Ã--- today '' not
because incomes fell, but because the credit ceiling rose faster than
incomes and prices followed it.

\begin{center}\rule{0.5\linewidth}{0.5pt}\end{center}

\subsubsection{2.5 The Integrated Feedback
Loop}\label{the-integrated-feedback-loop}

The four dimensions do not operate independently. They form a coupled
system with reinforcing feedback loops. The complete cascade '' from
current conditions to potential crisis '' develops as follows:

\begin{enumerate}
\def\labelenumi{\arabic{enumi}.}
\tightlist
\item
  \textbf{Affordability stress} (8``10Ã--- P/E) and turnover collapse
  (3.7\%) mean the market is already clearing by rationing.
\item
  As turnover falls further, \textbf{absorption tightens}: developers
  reduce volumes, supply falls toward 80,000``90,000 homes annually.
\item
  Fewer transactions mean \textbf{fewer mortgage originations}: M4
  growth decelerates from \textasciitilde3\% toward 1``2\%.
\item
  Slower money supply growth reduces \textbf{household purchasing
  power}: consumption weakens, business confidence falls.
\item
  \textbf{Unemployment rises} (1``2 quarter lag): wage growth stagnates,
  labour market slack builds.
\item
  Fixed debt on lower nominal income shifts the \textbf{DSR distribution
  rightward}: vulnerable cohorts move into elevated stress.
\item
  \textbf{Credit impairment spikes}: arrears accelerate from 0.9\%
  toward 1.5``2.0\%; bank write-off losses accumulate.
\item
  Banks tighten lending standards to protect capital: \textbf{mortgage
  approvals fall} further, from \textasciitilde65,000/month toward
  \textasciitilde45,000.
\item
  Accelerated decline in originations \textbf{contracts M4 further} ''
  potentially into negative territory.
\item
  \textbf{Deflationary spiral}: nominal GDP growth collapses; real debt
  burdens rise; employment deteriorates further.
\end{enumerate}

This cycle typically develops over 12``24 months from initial shock to
full crisis manifestation. Endogenous triggers are less likely '' demand
has already been substantially rationed by affordability. Exogenous
triggers are more probable: a geopolitical shock, a sharp unemployment
rise, rate volatility, or financial contagion from elsewhere. Any
trigger pushing unemployment above 5\% initiates the cascade described
above.

{\def\LTcaptype{none} % do not increment counter
\begin{longtable}[]{@{}
  >{\raggedright\arraybackslash}p{(\linewidth - 8\tabcolsep) * \real{0.2000}}
  >{\raggedright\arraybackslash}p{(\linewidth - 8\tabcolsep) * \real{0.2000}}
  >{\raggedright\arraybackslash}p{(\linewidth - 8\tabcolsep) * \real{0.2000}}
  >{\raggedright\arraybackslash}p{(\linewidth - 8\tabcolsep) * \real{0.2000}}
  >{\raggedright\arraybackslash}p{(\linewidth - 8\tabcolsep) * \real{0.2000}}@{}}
\toprule\noalign{}
\begin{minipage}[b]{\linewidth}\raggedright
\textbf{Phase}
\end{minipage} & \begin{minipage}[b]{\linewidth}\raggedright
\textbf{Months 0``3}
\end{minipage} & \begin{minipage}[b]{\linewidth}\raggedright
\textbf{Months 4``6}
\end{minipage} & \begin{minipage}[b]{\linewidth}\raggedright
\textbf{Months 7``12}
\end{minipage} & \begin{minipage}[b]{\linewidth}\raggedright
\textbf{Months 13``24}
\end{minipage} \\
\midrule\noalign{}
\endhead
\bottomrule\noalign{}
\endlastfoot
\textbf{Turnover} & 3.7\% & Falls to \textasciitilde3.0\% & Falls to
\textasciitilde2.5\% & Cliff risk \textless2\% \\
\textbf{M4 Growth} & \textasciitilde3\% & Decelerates to
\textasciitilde2.5\% & Falls to 1``2\% & Possible contraction \\
\textbf{Unemployment} & 3.8\% & 4.1``4.3\% & 4.5``5.0\% & 5.5``6.5\%+ \\
\textbf{Mortgage Arrears} & 0.9\% & Rising to \textasciitilde1.1\% &
Accelerating to 1.3``1.5\% & Crisis levels \textgreater2\% \\
\textbf{Credit Standards} & Tightening starts & Accelerated tightening &
Major tightening & Severe rationing \\
\end{longtable}
}

\textbf{Systemic risk assessment (March 2026):}

\begin{itemize}
\tightlist
\item
  \textbf{Short-term (0``12 months):} Stable; stress unlikely without
  exogenous shock.
\item
  \textbf{Medium-term (1``3 years):} Vulnerable; endogenous fragilities
  building.
\item
  \textbf{Long-term (3``5 years):} Unsustainable; crisis probable
  without structural reform.
\end{itemize}

Probability distribution: stable scenario 40``50\%; moderate stress
30``40\%; crisis scenario (\textgreater2\% arrears, financial stability
concerns) 15``25\%.

\begin{center}\rule{0.5\linewidth}{0.5pt}\end{center}

\protect\phantomsection\label{chapter-3}
\subsection{Chapter 3 '' Why Standard Responses Fail, and What
Structural Reform
Requires}\label{chapter-3-why-standard-responses-fail-and-what-structural-reform-requires}

The four-dimensional diagnosis in Chapter 2 has a direct policy
implication: if affordability deterioration is driven by credit-channel
dynamics, absorption-rate constraints, and monetary dependency '' rather
than by a simple physical shortage '' then standard policy responses
will not resolve it. They address symptoms while the structural engine
of exclusion continues to run. This chapter examines each standard
response category in turn, explains precisely why it fails, and sets out
what structural reform would actually require.

\begin{center}\rule{0.5\linewidth}{0.5pt}\end{center}

\subsubsection{3.1 Why Standard Responses
Fail}\label{why-standard-responses-fail}

\paragraph{Monetary Policy: Rate Cuts}\label{monetary-policy-rate-cuts}

Reducing the policy rate improves mortgage affordability and increases
mortgage demand '' typically by 10``20\% for a 100``150 basis point cut.
However, rate cuts cannot overcome the absorption rate constraint.
Developers maintain current pacing regardless of demand stimulus:
releasing more homes faster would depress the asset values that underpin
their balance sheets and land option portfolios. The net effect of
monetary easing with an unchanged supply ceiling is \textbf{price
inflation, not affordability improvement}. Verified demand rises;
effective demand buyers are still rationed out; prices are bid up to the
new, higher credit ceiling.

The historical evidence is unambiguous. Help to Buy (£28 billion
invested, 2013``2023) correlated with house prices rising 84\% while
supply rose only 16\%. Demand stimulus without supply-side structural
change redistributes purchasing power upward through the price level.

\paragraph{Planning Reform and Land
Release}\label{planning-reform-and-land-release}

Planning reform increases the number of plots on which development is
permitted. It does not increase the absorption rate within the outlet
system. Developers respond by holding more land in option agreements,
delaying development starts, and pacing build-out more slowly across a
larger portfolio. The number of sites with planning permission sitting
undeveloped has grown over the past decade even as approval rates
increased. Effective planning approvals have risen; housing completion
rates have remained stuck at or below 240,000 annually.

Planning reform is necessary but not sufficient. It addresses outlet
growth potential but cannot overcome developer pacing strategy if profit
incentives remain unchanged. The constraint is institutional and
incentive-driven, not regulatory.

\paragraph{Fiscal Supply Stimulus}\label{fiscal-supply-stimulus}

Direct government investment in housing could theoretically bypass the
private developer absorption ceiling. However, the scale required is
politically infeasible under current frameworks: a programme sufficient
to reach 300,000+ annual completions would require £40``60 billion
annually plus institutional development capacity that does not currently
exist. Even substantial programmes at achievable scale mainly substitute
for private supply rather than expanding total delivery. The absorption
ceiling shifts upward only if a genuinely different delivery model ''
with different incentive structures '' is deployed at scale.

\paragraph{Macroprudential Tweaks and Temporary
Guarantees}\label{macroprudential-tweaks-and-temporary-guarantees}

Adjusting loan-to-income flow limits, extending mortgage guarantee
schemes, or relaxing stress-test thresholds increases verified demand
without changing the structural conditions that generate the
verified-demand trap. As shown in Chapter 2, the gap between verified
demand and effective demand is the engine of long-run price inflation.
Widening access to high-LTV lending expands the cohort of borrowers
exposed to negative equity risk in the next downturn, increasing
systemic fragility while providing only temporary affordability relief.

The Basel III regulatory paradox compounds this: tighter capital
requirements applied to a mortgage-dominated system fall hardest on the
marginal borrower. Macroprudential tweaks that ease access at the margin
while leaving the structural concentration intact are, at best, symptom
management. At worst, they increase the amplitude of the next cycle.

\begin{center}\rule{0.5\linewidth}{0.5pt}\end{center}

\subsubsection{3.2 The Policy Table: Symptoms vs
Structure}\label{the-policy-table-symptoms-vs-structure}

{\def\LTcaptype{none} % do not increment counter
\begin{longtable}[]{@{}
  >{\raggedright\arraybackslash}p{(\linewidth - 4\tabcolsep) * \real{0.3333}}
  >{\raggedright\arraybackslash}p{(\linewidth - 4\tabcolsep) * \real{0.3333}}
  >{\raggedright\arraybackslash}p{(\linewidth - 4\tabcolsep) * \real{0.3333}}@{}}
\toprule\noalign{}
\begin{minipage}[b]{\linewidth}\raggedright
\textbf{Problem loop}
\end{minipage} & \begin{minipage}[b]{\linewidth}\raggedright
\textbf{Symptom-focused lever}
\end{minipage} & \begin{minipage}[b]{\linewidth}\raggedright
\textbf{Structural lever (examples)}
\end{minipage} \\
\midrule\noalign{}
\endhead
\bottomrule\noalign{}
\endlastfoot
Mortgage-credit dominance ties financial stability to housing cycles &
Macroprudential tweaks; temporary mortgage guarantees & Ring-fenced
affordable housing credit (New Circuit of Credit); diversification of
bank credit toward productive business lending \\
Thin-market clearing excludes mortgaged households & Stamp duty
holidays; first-time buyer grants; Help to Buy & Shared equity models;
cooperative homeownership; community land trusts \\
Developer pacing constrains near-term throughput & Planning speed
incentives; build-out rate requirements & Counter-cyclical public
delivery vehicles; cooperative and community-led development; public
build-to-live options \\
M4 dependency on housing credit creates monetary fragility &
Quantitative easing; rate cuts & Diversified credit creation through SME
and productive investment channels; New Circuit of Credit insulated from
asset-bidding dynamics \\
Verified demand inflation drives prices above effective demand & Income
multiple caps; stress-test adjustments & Land value capture; community
land trusts separating land from housing ownership; decommodification of
the land component \\
\end{longtable}
}

\begin{center}\rule{0.5\linewidth}{0.5pt}\end{center}

\subsubsection{3.3 What Structural Reform
Requires}\label{what-structural-reform-requires}

The regime diagnosis implies that resolution requires addressing all
four dimensions simultaneously. Structural reform that targets only one
dimension will be absorbed by the others. The following are the minimum
necessary components of a structural response.

\paragraph{A New Circuit of Credit}\label{a-new-circuit-of-credit}

The core monetary problem is that affordable housing delivery is
tethered to private-market credit cycles. When credit tightens,
affordable supply collapses; when credit loosens, it inflates prices
rather than expanding access. The Home@ix framework proposes a
\textbf{ring-fenced New Circuit of Credit (NCC)} '' a dedicated funding
channel for affordable housing delivery that does not compete with, or
bid up, existing stock. The NCC decouples affordable supply from
private-market freeze/thaw cycles and provides the stabiliser term in
the Affordable Housing Need identity:

\[AN = HD + (HM \cdot P \cdot AR \cdot D \cdot T \cdot PVC \cdot (F - 1)) + NCC\]

At baseline (\[F = 1\]), the middle term vanishes and need equals
effective demand plus the credit stabiliser. The NCC is the mechanism
that makes this identity solvable in practice '' without it, need always
exceeds supply when private credit contracts.

\paragraph{Decommodification: Community Land Trusts and Cooperative
Models}\label{decommodification-community-land-trusts-and-cooperative-models}

Separating land ownership (held collectively) from housing ownership
(held individually) breaks the speculative premium that drives the
verified-demand trap. Community land trusts have demonstrated the
capacity to reduce effective price-to-earnings ratios from 8``10Ã--- to
5``6Ã--- in comparable markets. Cooperative homeownership replaces
speculative investor ownership with participatory community control,
removing the incentive to pace sales for margin optimisation. Evidence
from Vienna (60\% of population in social housing, affordable, with
middle-class participation), Singapore (90\% homeownership via state-led
development), and North Dakota (community banking model with near-zero
housing default rates) demonstrates that alternative models can achieve
both affordability and access at scale.

\paragraph{Counter-Cyclical Delivery
Vehicles}\label{counter-cyclical-delivery-vehicles}

Overcoming the absorption rate constraint requires delivery vehicles
with different incentive structures '' organisations for which volume,
not margin, is the primary objective. Public development corporations,
housing associations operating at scale, and community-led developers
can pace delivery counter- cyclically: building more when private
developers pull back, and absorbing demand that the private market
cannot serve. This does not require the state to replace the private
market; it requires the state to fill the structural gap that the
private market's incentive structure creates.

\paragraph{Credit Concentration
Diversification}\label{credit-concentration-diversification}

Reducing the 60\% mortgage concentration in UK bank lending requires
active policy to redirect credit toward productive business investment
'' particularly SME lending, which has been structurally weak since the
GFC. This is partly a regulatory question (risk-weighting reform to
reduce the capital efficiency advantage of mortgage lending over
business lending) and partly an institutional question (building the
lending infrastructure for productive credit that was dismantled in the
post-GFC consolidation of the banking sector).

\begin{center}\rule{0.5\linewidth}{0.5pt}\end{center}

\subsubsection{3.4 The FAIR Indicator as a Reform
Monitor}\label{the-fair-indicator-as-a-reform-monitor}

Structural reform of the kind described above operates over years, not
quarters. During the transition, policymakers and researchers need a
forward-looking monitor that detects whether the regime is improving or
deteriorating '' and does so before the deterioration is visible in
conventional affordability statistics.

This is the operational role of \textbf{Home@ix FAIR}. By combining the
credit``price wedge (detecting decoupling between prices and the credit
base that finances them) with the turnover term (detecting thin-market
clearing), FAIR provides a quarterly signal of regime direction. A
sustained FAIR above 20 indicates that the structural conditions for
exclusion are intensifying, regardless of what headline price indices
show. A sustained FAIR below zero indicates that credit and market-depth
conditions are improving '' though, as the interpretation bands note,
very negative FAIR readings can themselves be stress-driven (post-crash
credit collapse rather than genuine improvement in access).

The formal derivation, reproducible construction, and backtesting
framework for FAIR are developed in the papers that follow. The chapters
above have established the diagnostic foundation: the regime that FAIR
is designed to detect, the mechanisms through which it operates, and the
structural reforms that would, if implemented, cause FAIR to sustainably
improve.

\textbf{The bottom line:} the UK housing system is not suffering from
scarcity. It is suffering from financialisation and structural
misallocation. With 60\% of bank credit concentrated in mortgages, M4
money supply vulnerable to housing shocks, developer pacing constraining
supply regardless of planning policy, and 35\% of transactions
controlled by cash buyers while mortgaged households are rationed out,
the system exhibits all the characteristics of an unsustainable regime.
Standard policy cannot break the feedback loops because it leaves the
incentive structures that generate them intact. Structural reform
addressing all four dimensions simultaneously is necessary. Without it,
crisis is probable within 2``5 years. With it, affordability, access,
and financial stability are achievable simultaneously.

\begin{center}\rule{0.5\linewidth}{0.5pt}\end{center}

\subsection{Bibliography}\label{bibliography}

Sources drawn upon across all parts of this working book. Where a formal
publisher or journal was not present on the source document, the
document type is noted for clarity.

\begin{enumerate}
\def\labelenumi{\arabic{enumi}.}
\tightlist
\item
  \protect\phantomsection\label{ref-archer2021}{\textbf{Archer, Tom, and
  Ian Cole.} ``The financialisation of housing production: exploring
  capital flows and value extraction among major housebuilders in the
  UK.'' \emph{Journal of Housing and the Built Environment} 36 (2021):
  1367``1387.}
\item
  \protect\phantomsection\label{ref-barker2004}{\textbf{Barker, Kate.}
  \emph{Barker Review of Housing Supply '' Securing our Future Housing
  Needs.} Report. HM Treasury, 2004.}
\item
  \protect\phantomsection\label{ref-davis1982}{\textbf{Davis, E. P., and
  I. D. Saville.} ``Mortgage lending and the housing market.''
  \emph{Bank of England Quarterly Bulletin}, Q3 (September 1982):
  390``398.}
\item
  \protect\phantomsection\label{ref-engdahl2012}{\textbf{Engdahl, F.
  William.} \emph{Myths, Lies and Oil Wars.} Wiesbaden: edition.engdahl,
  2012.}
\item
  \protect\phantomsection\label{ref-foye-forthcoming}{\textbf{Foye,
  Chris.} ``Framing the housing crisis: How think-tanks frame politics
  and science to advance policy agendas.'' \emph{Journal of Housing and
  the Built Environment.} Forthcoming.}
\item
  \protect\phantomsection\label{ref-lewis-opus}{\textbf{Lewis, Roger.}
  \emph{Affordable Housing OPUS, The Home@ix Formula.} Unpublished
  Draft/Paper.}
\item
  \protect\phantomsection\label{ref-lewis-linkedin-data}{\textbf{Lewis,
  Roger.} ``The Data Has Caught Up'' '' plain-English commentary
  version. \emph{LinkedIn Pulse}, 25 February 2026.}
\item
  \protect\phantomsection\label{ref-lewis-monetary}{\textbf{Lewis,
  Roger.} \emph{The Ending of the Long Monetary Expansion Cycle and a
  Brave New World of Housing Realism.} Unpublished Draft/Report.}
\item
  \protect\phantomsection\label{ref-lewis-bumper}{\textbf{Lewis, Roger.}
  \emph{The Home@ix Bumper Sticker.} Report/Statement.}
\item
  \protect\phantomsection\label{ref-lewis-framework}{\textbf{Lewis,
  Roger.} \emph{Home@ix Affordable Homes, A Framework of Understanding.}
  Report/White Paper.}
\item
  \protect\phantomsection\label{ref-lewis-linkedin-credit}{\textbf{Lewis,
  Roger.} ``Housing Credit, Broad Money, and the Real Economy: A UK View
  (or, When `Macroprudence' Becomes a Housing Scheme).'' \emph{LinkedIn
  Pulse}, 23 February 2026.}
\item
  \protect\phantomsection\label{ref-lewis-tensteps}{\textbf{Lewis,
  Roger.} \emph{Ten Steps to Affordable Housing: The New Circuit of
  Credit (The Circle of Blame, Volume 1).} Self-Published, 2025.}
\item
  \protect\phantomsection\label{ref-lewis-linkedin-freeze}{\textbf{Lewis,
  Roger.} ``When the Housing Market `Freezes': What Turnover, Cash
  Share, and Supply Mix Say About UK Housing Stress.'' \emph{LinkedIn
  Pulse}, 23 February 2026.}
\item
  \protect\phantomsection\label{ref-lietaer}{\textbf{Lietaer, Bernard,
  and Jacqui Dunne.} \emph{Rethinking Money: How New Currencies Turn
  Scarcity Into Prosperity.} Book.}
\item
  \protect\phantomsection\label{ref-meen2020}{\textbf{Meen, Geoffrey,
  and Christine Whitehead.} \emph{Understanding Affordability: The
  Economics of Housing Markets.} Bristol: Bristol University Press,
  2020.}
\item
  \protect\phantomsection\label{ref-mmcplaybook2024}{\textbf{MMC
  Playbook.} \emph{Social Rent Housing at Pace: The MMC Playbook. A
  Playbook for Local Authorities.} Report, May 2024.}
\item
  \protect\phantomsection\label{ref-mulheirn2019}{\textbf{Mulheirn,
  Ian.} \emph{Tackling the UK Housing Crisis: Is Supply the Answer?}
  Report. CaCHE (UK Collaborative Centre for Housing Evidence), August
  2019.}
\item
  \protect\phantomsection\label{ref-naess2020}{\textbf{Næss-Schmidt,
  Sigurd, Jonas Bjarke Jensen, Hendrik Ehmann, and Benjamin
  Christiansen.} \emph{Impact of the Final Basel III Framework in
  Sweden: Effects on the Banking Market and the Real Economy.}
  Copenhagen Economics for the Swedish Bankers' Association, February
  2020.}
\item
  \protect\phantomsection\label{ref-piddington2017}{\textbf{Piddington,
  Justine, Simon Nicol, Helen Garrett, and Matthew Custard.} \emph{The
  Housing Stock of The United Kingdom.} Report. BRE Trust, 2017.}
\item
  \protect\phantomsection\label{ref-savills2023}{\textbf{Savills
  Research.} \emph{A New Normal for Housebuilding? The Importance of
  Sales Outlets in a Market Without Help to Buy.} Report for Richborough
  Estates and LPDF, 1 March 2023.}
\item
  \protect\phantomsection\label{ref-sayre2020}{\textbf{Sayre, Brian.}
  \emph{Disrupt a Broken Industry '' The Industrial Construction
  Sandbox.} Report. Shadow Labs LLC, August 2020.}
\item
  \protect\phantomsection\label{ref-shin-bis}{\textbf{Shin, Hyun Song.}
  \emph{Macroprudential Policies Beyond Basel III.} BIS Papers No.~60.}
\item
  \protect\phantomsection\label{ref-shiller2015}{\textbf{Shiller, Robert
  J.} \emph{Irrational Exuberance} (Revised and Expanded Third Edition).
  Princeton: Princeton University Press, 2015.}
\item
  \protect\phantomsection\label{ref-starkey2018}{\textbf{Starkey,
  Maurice.} ``Utilising the Quantity Theory of Credit to Understand the
  Causes of the 2007 Financial Crisis.'' Economics Network / Creative
  Commons, 2018.}
\item
  \protect\phantomsection\label{ref-sfrg2010}{\textbf{Systemic Fiscal
  Reform Group.} ``The Principles of Tax Policy.'' Written evidence
  submitted to a Parliamentary committee, Session 2010``11.}
\item
  \protect\phantomsection\label{ref-whitehead2020}{\textbf{Whitehead,
  Christine, and Peter Williams.} \emph{Thinking Outside the Box:
  Exploring Innovations in Affordable Home Ownership.} Report. London
  School of Economics / Housing Evidence, 11 November 2020.}
\end{enumerate}

\protect\phantomsection\label{part1}
\subsection{Part I '' The 2024 Foundations: Market Function Before
Price}\label{part-i-the-2024-foundations-market-function-before-price}

Home@ix began from an observation that sounds simple but changes the
whole measurement problem: in UK housing stress regimes, the market does
not reliably clear through orderly repricing. It often clears through
\textbf{rationing participation}. Turnover collapses, transaction chains
fail, mortgage-dependent households are screened out, and the observed
price index becomes a \emph{selected statistic} produced by a thinner,
wealthier transacting set.

Part I sets the domain terms for the rest of the book. Before we define
FAIR (the derived, reproducible regime indicator), we anchor the
foundational claim with a canonical diagnostic module: turnover (market
liquidity), cash-versus-mortgage dynamics (access to finance), and the
balance between new build and old stock activity (delivery and churn
capacity).

\textbf{Working proposition:} affordability is not only a burden ratio;
it is an \textbf{access-and-allocation regime outcome}. In thin markets,
price signals can lag while exclusion and mobility failure worsen in
real time.

\begin{center}\rule{0.5\linewidth}{0.5pt}\end{center}

\protect\phantomsection\label{part1-exhibits}
\subsubsection{1. A canonical diagnostic module (market depth, access,
and
mix)}\label{a-canonical-diagnostic-module-market-depth-access-and-mix}

The exhibits below are treated as a single system lens. Read together,
they diagnose whether the housing system is clearing via price, or via
\emph{quantity and composition} (thin-market clearing).

\begin{center}\rule{0.5\linewidth}{0.5pt}\end{center}

\protect\phantomsection\label{part1-fig1}
\subsubsection{Figure 1 '' Turnover and cash share
(annual)}\label{figure-1-turnover-and-cash-share-annual}

Turnover is the housing market™s \textbf{functionality / liquidity}
indicator: how many homes change hands relative to the stock. Cash share
is a \textbf{participation / exclusion} indicator: when cash share rises
alongside weak turnover, it is consistent with mortgage-dependent
households being rationed out by rates, underwriting, or deposit
constraints.

Figure 1. Stock turnover and cash share (annual).

\begin{center}\rule{0.5\linewidth}{0.5pt}\end{center}

\protect\phantomsection\label{part1-fig2a}
\subsubsection{Figure 2a '' Transaction levels: old stock vs new build
(annual)}\label{figure-2a-transaction-levels-old-stock-vs-new-build-annual}

œHousing activity is not one market. Old stock dominates volume and is
more \textbf{credit- and chain-sensitive}. New build volumes reflect
additional constraints (planning, labour, build finance, developer
pacing). The split helps distinguish \emph{access/finance stress} from
\emph{delivery/capacity stress}.

Figure 2a. Transaction levels: old stock vs new build (annual).

\begin{center}\rule{0.5\linewidth}{0.5pt}\end{center}

\protect\phantomsection\label{part1-fig2b}
\subsubsection{Figure 2b '' Transaction mix shares
(annual)}\label{figure-2b-transaction-mix-shares-annual}

Mix shares reveal \textbf{compositional clearing}: who is still able to
transact when the market thins. Shifts in shares can signal allocation
stress even when headline prices appear stable.

Figure 2b. Transaction mix shares (annual).

\begin{center}\rule{0.5\linewidth}{0.5pt}\end{center}

\protect\phantomsection\label{part1-fig3}
\subsubsection{Figure 3 '' Mortgage outstanding composition by lender
type
(annual)}\label{figure-3-mortgage-outstanding-composition-by-lender-type-annual}

This exhibit shows how the stock of mortgage lending is distributed
across lender types. Shifts in composition matter because œcredit
availability is not a single dial: it is mediated by lender balance
sheets, underwriting posture, and the kinds of borrowers each segment
serves.

Figure 3. Mortgage outstanding composition by lender type (annual, share
of total).

\begin{center}\rule{0.5\linewidth}{0.5pt}\end{center}

\protect\phantomsection\label{homeix-identity}
\subsection{Home@ix '' The Identity (as developed) and the Reproducible
Build}\label{homeix-the-identity-as-developed-and-the-reproducible-build}

Home@ix uses accounting identities to describe the housing system as a
throughput machine. The point is not œmath for effect. The point is
discipline: when the system is constrained by absorption, credit gates,
delivery friction, and value-capture reality, output behaves like a set
of multipliers. If one dial is near zero, the product collapses.

\protect\phantomsection\label{homeix-identity-ams}
\subsubsection{Identity 1: Affordable Market Supply
(AMS)}\label{identity-1-affordable-market-supply-ams}

We start with the supply identity: the affordable supply the system can
deliver, given real-world throttles.

\textbf{Affordable Market Supply}:

\[AMS = (HM \cdot P \cdot AR \cdot D \cdot T \cdot PVC) + HS\]

\begin{itemize}
\tightlist
\item
  \textbf{HM}: housebuilding market capacity (deliverable output, not
  theoretical land)
\item
  \textbf{P}: affordable proportion of delivery
\item
  \textbf{AR}: absorption rate (how fast homes can be sold/let without
  destabilising price)
\item
  \textbf{D}: diversity (tenure/product mix widening take-up)
\item
  \textbf{T}: throughput (permission †' start †' completion conversion
  efficiency)
\item
  \textbf{PVC}: planning value capture that works in practice
\item
  \textbf{HS}: existing affordable housing stock
\end{itemize}

\protect\phantomsection\label{homeix-identity-an}
\subsubsection{Identity 2: Affordable Housing Need (AN),
derived}\label{identity-2-affordable-housing-need-an-derived}

Need persists (and often rises) when supply is tethered to
private-market throughput and absorption constraints. Home@ix expresses
this by introducing a sequencing/fast-tracking factor \textbf{F} and a
stabiliser term: a ring-fenced \textbf{New Circuit of Credit}
(\textbf{NCC}) that funds affordable delivery without simply bidding up
existing stock.

\textbf{Start from the expanded form}:

\[AN = HD + (HM \cdot P \cdot AR \cdot D \cdot T \cdot PVC \cdot F) + HS - AMS + NCC\]

\textbf{Substitute}
\[AMS = (HM \cdot P \cdot AR \cdot D \cdot T \cdot PVC) + HS\]
\textbf{and simplify}:

\[AN = HD + (HM \cdot P \cdot AR \cdot D \cdot T \cdot PVC \cdot (F - 1)) + NCC\]

\begin{itemize}
\tightlist
\item
  \textbf{HD}: HomeMaker effective demand (need in households, not only
  œbankable demand this quarter)
\item
  \textbf{F}: sequencing/fast-tracking factor (the sign convention is:
  baseline at \[F=1\])
\item
  \textbf{NCC}: New Circuit of Credit Creation (the stabiliser that
  decouples affordable supply from private-market freeze/thaw cycles)
\end{itemize}

\protect\phantomsection\label{repo-structure}
\subsubsection{Reproducibility: from canonical inputs to the figures in
this
book}\label{reproducibility-from-canonical-inputs-to-the-figures-in-this-book}

This repository is organised so a reader can move from source data to
computed series to plotted exhibits''and regenerate the print bundle.
The print folder you are reading is a portable snapshot of that
pipeline.

\begin{itemize}
\tightlist
\item
  \textbf{Canonical inputs (source of truth):}
  \texttt{inputs/canonical/}
\item
  \textbf{Intermediate processing (rebuildable cache):}
  \texttt{build/processed/}
\item
  \textbf{Generated artefacts referenced by papers:}
  \texttt{outputs/figures/}, \texttt{outputs/fair\_assets/},
  \texttt{outputs/draft\_paper\_assets/}
\item
  \textbf{Portable publication bundle:} \texttt{publication/print\_*/}
  containing \texttt{index.html} and copied \texttt{assets/}
\end{itemize}

Operationally, the maths œcomputes as deterministic column derivations
and transforms inside the Python scripts, and the figures are the audit
trail of those computations.

\begin{itemize}
\tightlist
\item
  \texttt{sweep\_up/inbox/homeatix\_annual\_exhibits.py} '' reads
  canonical data and generates the Part I exhibits into
  \texttt{assets/figures/}
\item
  \texttt{sweep\_up/inbox/assemble\_full\_ebook.py} '' writes the
  combined \texttt{index.html} into the print bundle and normalises
  asset paths to \texttt{assets/}
\end{itemize}

\begin{center}\rule{0.5\linewidth}{0.5pt}\end{center}

\protect\phantomsection\label{part1-fig4}
\subsubsection{Figure 4 '' Indexed price vs indexed turnover
(annual)}\label{figure-4-indexed-price-vs-indexed-turnover-annual}

When turnover collapses, price indices can lag or understate stress
because fewer transactions generate weaker price discovery. Liquidity
risk becomes macro-relevant.

Figure 4. Indexed price vs turnover (annual).

\begin{center}\rule{0.5\linewidth}{0.5pt}\end{center}

\protect\phantomsection\label{part1-cashshare-ew}
\subsubsection{Supplement '' Cash purchases share (England vs Wales,
quarterly)}\label{supplement-cash-purchases-share-england-vs-wales-quarterly}

A cross-nation comparison helps separate broad regime shifts from local
noise. Divergences can reflect geography-specific credit conditions,
investor mix, or compositional effects.

Supplement. Cash purchases as a proportion of total transactions
(quarterly): England vs Wales.

\begin{center}\rule{0.5\linewidth}{0.5pt}\end{center}

\protect\phantomsection\label{part1-bridge-to-identity}
\subsubsection{2. Why this forces an accounting identity (and later,
FAIR)}\label{why-this-forces-an-accounting-identity-and-later-fair}

If the system clears through \textbf{quantities and composition} under
stress, then an affordability framework must track access-to-finance and
market depth/throughput '' not just price levels.

\textbf{Bridge statement:} Part I defines the failure mode (thin-market
clearing). FAIR later converts that diagnosis into a reproducible
quarterly monitor.

\begin{center}\rule{0.5\linewidth}{0.5pt}\end{center}

\protect\phantomsection\label{part1-source-note}
\subsubsection{3. Source note
(reproducibility)}\label{source-note-reproducibility}

\begin{itemize}
\tightlist
\item
  \textbf{Exhibits generator:}
  \texttt{sweep\_up/inbox/homeatix\_annual\_exhibits.py}
\item
  \textbf{Bundled figures:} \texttt{assets/figures/} (relative to this
  \texttt{index.html})
\end{itemize}

\begin{center}\rule{0.5\linewidth}{0.5pt}\end{center}

\protect\phantomsection\label{followup}
\subsection{Follow-up Paper}\label{follow-up-paper}

\begin{center}\rule{0.5\linewidth}{0.5pt}\end{center}

\section{\texorpdfstring{\textbf{Home@ix Follow-up:} Mortgage-Credit
Concentration, M4 Dependency, and UK Housing Affordability Regimes ''
With a Reproducible Early-Warning Indicator
(FAIR)}{Home@ix Follow-up: Mortgage-Credit Concentration, M4 Dependency, and UK Housing Affordability Regimes '' With a Reproducible Early-Warning Indicator (FAIR)}}\label{homeix-follow-up-mortgage-credit-concentration-m4-dependency-and-uk-housing-affordability-regimes-with-a-reproducible-early-warning-indicator-fair}

\textbf{Author:} Home@ix (Working Paper / Draft)

\textbf{Date:} 25 February 2026

\subsection{Abstract}\label{abstract}

The UK housing affordability crisis is commonly presented as a
price-to-income problem. That framing is incomplete in regimes where
housing markets do not clear through price adjustment. Under stress,
clearing occurs through quantities and composition: turnover falls,
chains fail, mortgage-dependent households are rationed out, and the
observed price level becomes a lagging and selected statistic.

This paper connects four reinforcing dimensions: (i) credit impairment
risk in commercial banks under mortgage-credit concentration, (ii)
mortgage lending over-reliance and collateral feedback, (iii) structural
dependency of broad money growth on mortgage credit creation, and (iv)
affordability deterioration as an access-and-allocation outcome. It then
operationalises the diagnosis using \textbf{Home@ix FAIR}, a transparent
quarterly indicator designed as a \textbf{2``3 year forward regime
signal} for affordability understood as throughput and access. FAIR
combines a credit``price decoupling wedge with a market-depth (turnover)
term, optionally augmented by composition proxies.

\subsection{Executive Summary}\label{executive-summary}

UK housing stress is not only about œhigh prices. It is also about
\textbf{who can transact} and \textbf{how the market clears}. When
mortgage credit tightens or uncertainty rises, the UK market often
clears via \textbf{reduced participation} rather than orderly repricing:
transactions fall, turnover collapses, and cash-rich buyers gain share.

\textbf{Core claim:} mortgage-credit concentration in banks links
financial stability and affordability. When mortgage flows weaken,
deposit creation weakens; when deposit growth weakens, the real economy
tightens; when the economy tightens, arrears risks rise; and when
arrears risks rise, credit supply tightens further. Affordability
deteriorates via allocation and throughput failure.

\subsubsection{Key findings (sense-checked, definitions required in
final)}\label{key-findings-sense-checked-definitions-required-in-final}

\begin{itemize}
\tightlist
\item
  \textbf{Mortgage concentration is systemic:} residential mortgages
  comprise a large share of UK bank private-sector lending stock,
  increasing transmission from housing shocks to bank balance sheets.
\item
  \textbf{Broad money is mortgage-sensitive:} when mortgage lending
  slows, broad money growth can undershoot nominal activity trends,
  reinforcing tight conditions.
\item
  \textbf{Near-term supply is throughput-constrained:} major
  housebuilders pace sales within a narrow absorption band (sales per
  outlet per week), limiting market clearing via volume.
\item
  \textbf{Affordability is an allocation problem in stress:} turnover
  collapses and buyer composition shifts (cash share rises), so the
  market œclears by rationing access rather than by repricing alone.
\item
  \textbf{FAIR provides an operational monitor:} FAIR detects
  credit``price decoupling and thinning market depth, giving a
  forward-looking regime signal consistent with the diagnosis.
\end{itemize}

\subsection{1. The Integrated Mechanism: Why These Four Dimensions Move
Together}\label{the-integrated-mechanism-why-these-four-dimensions-move-together}

The four dimensions''bank credit quality, mortgage concentration, broad
money dynamics, and affordability''are often analysed separately. In
practice, they form a coupled system with feedback loops.

\subsubsection{1.1 System sketch (feedback
loops)}\label{system-sketch-feedback-loops}

\begin{itemize}
\tightlist
\item
  \textbf{Mortgage concentration loop:} house prices and employment
  shocks affect arrears and loss expectations, which tighten
  underwriting and reduce credit availability.
\item
  \textbf{Endogenous money loop:} mortgage lending creates deposits;
  weakening mortgage flows reduce broad money growth, which can tighten
  spending conditions and worsen employment.
\item
  \textbf{Thin-market loop:} tighter credit reduces the set of eligible
  buyers; turnover falls; observed prices become œstickier and more
  selected; access worsens even if the price index moves slowly.
\item
  \textbf{Developer pacing loop:} when demand weakens, builders protect
  margins by controlling release pace; volume cannot rise enough to
  clear via supply in the near term.
\end{itemize}

This paper treats œaffordability as the regime outcome of these loops,
not as a single ratio.

\subsection{2. Credit Impairment: Surface Health vs Structural
Fragility}\label{credit-impairment-surface-health-vs-structural-fragility}

Headline arrears and delinquency can look stable even as systemic
vulnerability rises, because compositional selection improves the
observed pool of borrowers while stress accumulates in marginal cohorts
(recent high-LTV buyers, BTL, renters-in-distress feeding landlord
stress).

\subsubsection{2.1 Why low arrears can be
misleading}\label{why-low-arrears-can-be-misleading}

\begin{itemize}
\tightlist
\item
  \textbf{Selection effects:} improved observed arrears may partly
  reflect that the stressed tail stops transacting or is screened out of
  new lending.
\item
  \textbf{Threshold reporting:} œserious arrears definitions can miss
  earlier-stage stress that predicts later acceleration.
\item
  \textbf{BTL as leading indicator:} landlord cashflows are more
  directly exposed to rate changes and tenant stress; BTL arrears can
  precede broader deterioration.
\end{itemize}

\subsubsection{2.2 Debt servicing: risk rises
continuously}\label{debt-servicing-risk-rises-continuously}

If arrears probability increases roughly linearly beyond a low-stress
floor, then stability depends on distribution shifts, not just tail
events. Small rightward shifts in the debt-servicing distribution can
move a meaningful share of households into elevated stress.

\subsection{3. Mortgage Lending Concentration and M4
Dependency}\label{mortgage-lending-concentration-and-m4-dependency}

In an endogenous money system, bank lending and deposit creation are
linked. If mortgage lending dominates bank balance sheets, then
housing-credit cycles imprint onto broad money dynamics and can
propagate to the real economy.

\subsubsection{3.1 What œM4 dependency means
operationally}\label{what-ux153m4-dependency-means-operationally}

\begin{itemize}
\tightlist
\item
  \textbf{Mechanism:} new mortgage lending creates deposits; weak
  mortgage origination can weaken deposit growth.
\item
  \textbf{Interpretation:} slowing broad money growth is not
  mechanically œbad, but it becomes a stress signal when it coincides
  with tight credit, weak turnover, and rising delinquency risk.
\item
  \textbf{Empirical caution:} velocity and portfolio shifts matter; M4
  is best used as a regime indicator alongside housing-credit and
  turnover measures.
\end{itemize}

\subsection{4. Housing Affordability as Market Clearing: Access,
Allocation,
Throughput}\label{housing-affordability-as-market-clearing-access-allocation-throughput}

In stressed regimes, housing markets can fail to clear via prices and
instead clear via participation constraints: fewer mortgage-dependent
buyers qualify, chains break more often, and turnover collapses.

\subsubsection{4.1 Compositional clearing}\label{compositional-clearing}

\begin{itemize}
\tightlist
\item
  \textbf{Turnover collapse:} when transactions fall relative to stock,
  the market becomes thin; observed prices reflect a narrower set of
  buyers and properties.
\item
  \textbf{Buyer-type rationing:} higher cash shares can coexist with
  sticky prices because cash demand clears the marginal flow even as
  mortgaged households are excluded.
\item
  \textbf{Distributional harm:} renters and younger cohorts experience
  the binding constraint (access), even when average ratios look stable.
\end{itemize}

\subsubsection{4.2 The absorption-rate constraint (near
term)}\label{the-absorption-rate-constraint-near-term}

New-build supply often cannot surge quickly enough to restore clearing
through volume because developers pace sales (per outlet per week) to
optimise margins and manage risk. This creates a throughput constraint
that is institutional/incentive-driven, not purely physical.

\subsection{5. From Diagnosis to Tooling: The Home@ix FAIR
Indicator}\label{from-diagnosis-to-tooling-the-homeix-fair-indicator}

The diagnosis implies two observable precursors to allocation stress
that are available in aggregate data: \textbf{credit``price decoupling}
and \textbf{thin-market clearing}. FAIR is a reduced-form index that
combines both.

\subsubsection{5.1 Inputs (quarterly)}\label{inputs-quarterly}

\begin{itemize}
\tightlist
\item
  House prices (level): \texttt{avg\_house\_price\_gbp}
\item
  Mortgage stock (level, £m): \texttt{mb\_total\_gbp\_m}
\item
  Turnover (market depth): \texttt{turnover\_pct\_q}
\item
  Optional: new-build share: \texttt{newbuild\_share\_of\_transactions}
\item
  Quarter index: \texttt{period}; geography: \texttt{geo} (e.g., EW)
\end{itemize}

\subsubsection{5.2 Derived series}\label{derived-series}

Define year-on-year growth rates:

\[ g_t^P = YoY(P_t) \]

\[ g_t^{MB} = YoY(MB_t) \]

\[ g_t^{TO} = YoY(TO_t) \]

Define the credit``price wedge:

\[ W_t = g_t^P - g_t^{MB} \]

Optional new-build share change (YoY):

\[ \Delta NB_t = NB_t - NB\_{t-4} \]

\subsection{6. Baseline Normalisation (Core
Proviso)}\label{baseline-normalisation-core-proviso}

To avoid defining œnormal using crisis regimes and structural breaks,
z-scores are computed from pooled baseline windows:

\begin{itemize}
\tightlist
\item
  \textbf{Baseline windows:} 2003Q1``2007Q4 and 2013Q1``2019Q4
\item
  \textbf{Excluded:} 2008``2012 (GFC and aftermath) and 2020+ (Covid and
  subsequent regime breaks)
\end{itemize}

For any series \[X_t\], define:

\[ z(X_t) = \frac{X_t - \mu_X}{\sigma_X} \]

where \[\mu_X\] and \[\sigma_X\] are computed using baseline quarters
only.

\subsection{7. FAIR Definition, Interpretation, and
Monitoring}\label{fair-definition-interpretation-and-monitoring}

FAIR is scaled for interpretability and designed to be auditable.
Weights are intentionally simple and can be stress-tested.

\subsubsection{7.1 FAIR definition}\label{fair-definition}

Default (three-component) specification:

\[ FAIR_t = 100 \cdot \Big( 0.55 \cdot z(W_t) - 0.35 \cdot z(g_t^{TO}) + 0.10 \cdot z(\Delta NB_t) \Big) \]

If the new-build term is not available:

\[ FAIR_t = 100 \cdot \Big( 0.55 \cdot z(W_t) - 0.35 \cdot z(g_t^{TO}) \Big) \]

\subsubsection{7.2 Direction-of-flow}\label{direction-of-flow}

\[ \Delta FAIR_t = FAIR_t - FAIR\_{t-1} \]

The level indicates regime pressure; the first difference indicates
acceleration or easing.

\subsubsection{7.3 Practical interpretation
bands}\label{practical-interpretation-bands}

\begin{itemize}
\tightlist
\item
  \[FAIR \ge 50\]: strong deterioration regime risk
\item
  \[20 \le FAIR < 50\]: mild deterioration bias
\item
  \[-20 \le FAIR < 20\]: neutral/noisy
\item
  \[-50 \le FAIR < -20\]: mild improvement bias
\item
  \[FAIR < -50\]: strong improvement regime (often stress-driven)
\end{itemize}

\subsection{8. Evaluation: Event-Based Backtesting (Crash-Start
Definitions)}\label{evaluation-event-based-backtesting-crash-start-definitions}

FAIR can be evaluated as an early-warning tool using objective event
definitions rather than narrative dating.

\subsubsection{8.1 Crash-start events
(rule-based)}\label{crash-start-events-rule-based}

A œcrash start may be defined as a local price peak followed by a
drawdown of at least \[X\\] within \[Y\] quarters, with a minimum
separation (cooldown) of \[C\] quarters to avoid double-counting
clustered peaks.

\subsubsection{8.2 Illustrative warning
rules}\label{illustrative-warning-rules}

\begin{itemize}
\tightlist
\item
  A: sustained level stress: \[FAIR > 20\] for 2 quarters
\item
  B: sustained acceleration: \[\Delta FAIR > 5\] for 2 quarters
\item
  C: broad stress-worsening: \[FAIR > 0\] and \[\Delta FAIR \ge 0\]
\end{itemize}

\subsubsection{8.3 Metrics}\label{metrics}

\begin{itemize}
\tightlist
\item
  Lead time (quarters) from the most recent signal to each crash start
\item
  False-positive proxy: share of signal quarters not followed by any
  crash start within a forward window (e.g., 8 quarters)
\end{itemize}

\subsection{9. Outputs and Visualisations
(Reproducible)}\label{outputs-and-visualisations-reproducible}

This draft renders the latest FAIR assets directly from the Home@ix
output folders. All paths below are \textbf{relative to the server root}
(the folder you run \texttt{py\ -m\ http.server} in).

\subsubsection{9.1 FAIR headline level}\label{fair-headline-level}

Figure 1. Home@ix FAIR (level). Source:
assets/figures/fig\_fair\_level.png

Figure 1b. Home@ix FAIR (level). Source:
assets/fair\_assets/fig\_fair\_level.png

\begin{center}\rule{0.5\linewidth}{0.5pt}\end{center}

\subsubsection{9.2 Component
contributions}\label{component-contributions}

Figure 2. FAIR component contributions. Source:
assets/figures/fig\_fair\_contrib.png

Figure 2b. FAIR component contributions. Source:
assets/fair\_assets/fig\_fair\_contributions.png

\begin{center}\rule{0.5\linewidth}{0.5pt}\end{center}

\subsubsection{9.3 Direction-of-flow (animated
GIFs)}\label{direction-of-flow-animated-gifs}

Figure 3. Direction-of-flow animation. Source:
assets/figures/anim\_fair\_flow.gif

Figure 3b. Direction-of-flow animation. Source:
assets/fair\_assets/anim\_fair\_direction\_of\_flow.gif

\begin{center}\rule{0.5\linewidth}{0.5pt}\end{center}

\subsubsection{9.4 Backtest / evaluation
figures}\label{backtest-evaluation-figures}

Figure 4. Prices and FAIR with crash-start annotations. Source:
assets/draft\_paper\_assets/fig\_price\_and\_fair\_with\_crisis\_starts.png

Figure 5. Average lead time by signal rule. Source:
assets/draft\_paper\_assets/fig\_avg\_leadtime\_by\_signal.png

If any image fails to load, your Python server console will log a 404
showing the exact path requested.

\subsection{10. Policy and Structural Reform (Implications of the Regime
View)}\label{policy-and-structural-reform-implications-of-the-regime-view}

Standard levers (rate cuts, planning reform, time-limited subsidies) can
influence symptoms without breaking the feedback loops. If affordability
deterioration is driven by allocation and throughput failure under
mortgage-credit dominance, then structural reform must address credit
channel design and tenure/ownership pathways.

{\def\LTcaptype{none} % do not increment counter
\begin{longtable}[]{@{}
  >{\raggedright\arraybackslash}p{(\linewidth - 4\tabcolsep) * \real{0.3333}}
  >{\raggedright\arraybackslash}p{(\linewidth - 4\tabcolsep) * \real{0.3333}}
  >{\raggedright\arraybackslash}p{(\linewidth - 4\tabcolsep) * \real{0.3333}}@{}}
\toprule\noalign{}
\begin{minipage}[b]{\linewidth}\raggedright
\textbf{Problem loop}
\end{minipage} & \begin{minipage}[b]{\linewidth}\raggedright
\textbf{Symptom-focused lever}
\end{minipage} & \begin{minipage}[b]{\linewidth}\raggedright
\textbf{Structural lever (examples)}
\end{minipage} \\
\midrule\noalign{}
\endhead
\bottomrule\noalign{}
\endlastfoot
Mortgage-credit dominance ties stability to housing & Macroprudential
tweaks, temporary guarantees & Ring-fenced affordable housing credit;
alternative tenure finance \\
Thin-market clearing excludes mortgaged households & Stamp duty
holidays, small buyer grants & Shared equity/co-op models; community
land trusts \\
Developer pacing constrains near-term throughput & Planning speed,
incentives to build & Counter-cyclical delivery vehicles; public options
for build-to-live \\
\end{longtable}
}

\subsection{11. Limitations and
Extensions}\label{limitations-and-extensions}

FAIR is intentionally reduced-form. It captures two dominant mechanisms
but does not directly encode: arrears microdata, cash-buyer share,
investor/renter distress transmission, or explicit developer pacing
series. These are natural extensions that can be layered as satellite
indicators.

\subsection{12. Conclusion}\label{conclusion}

UK housing affordability is not only a price level story; it is a regime
story about who can transact and how the market clears when credit
tightens. Mortgage-credit concentration links affordability outcomes to
financial stability dynamics, and broad money conditions can amplify
stress. Home@ix FAIR translates this diagnosis into an operational,
reproducible monitor of regime risk by combining credit``price
decoupling with market depth.

\subsection{Reproducibility Note}\label{reproducibility-note}

Expected output directory contents (adjust to your workflow):

\begin{itemize}
\tightlist
\item
  \texttt{assets/fair\_assets/fair\_quarterly\_audit.csv}
\item
  \texttt{assets/figures/fig\_fair\_level.png}
\item
  \texttt{assets/figures/fig\_fair\_contrib.png}
\item
  \texttt{assets/figures/anim\_fair\_flow.gif}
\item
  \texttt{assets/fair\_assets/fig\_fair\_level.png}
\item
  \texttt{assets/fair\_assets/fig\_fair\_contributions.png}
\item
  \texttt{assets/fair\_assets/anim\_fair\_direction\_of\_flow.gif}
\item
  \texttt{assets/draft\_paper\_assets/fig\_price\_and\_fair\_with\_crisis\_starts.png}
\item
  \texttt{assets/draft\_paper\_assets/fig\_avg\_leadtime\_by\_signal.png}
\end{itemize}

\begin{center}\rule{0.5\linewidth}{0.5pt}\end{center}

\protect\phantomsection\label{fair-preface}
\subsection{FAIR '' Professional summary
(LinkedIn-style)}\label{fair-professional-summary-linkedin-style}

This section is a distilled, non-technical bridge into the FAIR working
paper. It is written for housing policy, finance, and economics
audiences.

\textbf{Thesis:} The UK housing crisis isn™t only about prices. In
stressed regimes, markets can œclear through \textbf{exclusion}, not
repricing.

When turnover collapses and mortgage-dependent households are screened
out, the reported price index becomes a \emph{selected statistic}
generated by a thinner, wealthier transacting set.

\subsubsection{The coupled system (four
loops)}\label{the-coupled-system-four-loops}

\begin{enumerate}
\def\labelenumi{\arabic{enumi}.}
\tightlist
\item
  \textbf{Bank credit quality and mortgage concentration}
\item
  \textbf{Mortgage lending and broad money dynamics} (deposit creation /
  M4 channel)
\item
  \textbf{Supply delivery constraints} and developer pacing within
  absorption bands
\item
  \textbf{Affordability} understood as access-and-allocation, not only
  burden ratios
\end{enumerate}

\subsubsection{Why œlow arrears can be a false
comfort}\label{why-ux153low-arrears-can-be-a-false-comfort}

Headline arrears can look stable even as vulnerability rises, because
the observed borrower pool can improve via \textbf{compositional
selection} (marginal cohorts stop transacting or are screened out of new
lending). Buy-to-let arrears can behave as a leading indicator when
landlord cashflows tighten.

\subsubsection{Introducing FAIR (a forward regime
signal)}\label{introducing-fair-a-forward-regime-signal}

FAIR is a quarterly indicator designed as a \textbf{2``3 year forward
regime signal} for affordability as access/allocation. It combines:

\begin{itemize}
\tightlist
\item
  \textbf{Credit``price decoupling} (a credit``price wedge)
\item
  \textbf{Market depth} (turnover dynamics relative to stock)
\item
  \emph{Optional:} new-build share dynamics as a composition proxy
\end{itemize}

\textbf{Transparent form (as used in this project):}

\[FAIR = 100 \times (0.55 \times \text{credit“price wedge} - 0.35 \times \text{turnover growth} + 0.10 \times \Delta\text{new-build share})\]

Components are z-scored against pooled baseline windows (2003``2007 and
2013``2019) to avoid structural-break periods.

\subsubsection{Interpretation bands}\label{interpretation-bands}

{\def\LTcaptype{none} % do not increment counter
\begin{longtable}[]{@{}ll@{}}
\toprule\noalign{}
FAIR level & Regime signal \\
\midrule\noalign{}
\endhead
\bottomrule\noalign{}
\endlastfoot
‰¥ 50 & Strong deterioration risk \\
20 to 50 & Mild deterioration bias \\
ˆ'20 to 20 & Neutral / noisy \\
ˆ'50 to ˆ'20 & Mild improvement bias \\
\textless{} ˆ'50 & Strong improvement (often stress-driven) \\
\end{longtable}
}

\subsubsection{Policy implication (what gets
missed)}\label{policy-implication-what-gets-missed}

If affordability deterioration is driven by \textbf{allocation and
throughput failure under mortgage-credit dominance}, then symptom levers
alone (temporary subsidies, short-run tax changes) can lag the regime
shift.

{\def\LTcaptype{none} % do not increment counter
\begin{longtable}[]{@{}
  >{\raggedright\arraybackslash}p{(\linewidth - 4\tabcolsep) * \real{0.3333}}
  >{\raggedright\arraybackslash}p{(\linewidth - 4\tabcolsep) * \real{0.3333}}
  >{\raggedright\arraybackslash}p{(\linewidth - 4\tabcolsep) * \real{0.3333}}@{}}
\toprule\noalign{}
\begin{minipage}[b]{\linewidth}\raggedright
Problem loop
\end{minipage} & \begin{minipage}[b]{\linewidth}\raggedright
Symptom lever
\end{minipage} & \begin{minipage}[b]{\linewidth}\raggedright
Structural lever
\end{minipage} \\
\midrule\noalign{}
\endhead
\bottomrule\noalign{}
\endlastfoot
Mortgage-credit dominance ties stability to housing & Macroprudential
tweaks, temporary guarantees & Ring-fenced affordable housing credit;
alternative tenure finance \\
Thin-market clearing excludes mortgaged households & Buyer grants,
transaction stimulus & Shared equity, co-op models, community land
trusts \\
Developer pacing constrains near-term throughput & Planning speed
incentives & Counter-cyclical delivery vehicles; public build-to-live
options \\
\end{longtable}
}

\textbf{Bridge:} The FAIR paper that follows formalises the measurement,
definitions, and reproducible build for this indicator.

\begin{center}\rule{0.5\linewidth}{0.5pt}\end{center}

\begin{center}\rule{0.5\linewidth}{0.5pt}\end{center}

\section{Home@ix FAIR '' The Maths
Explained}\label{homeix-fair-the-maths-explained}

This paper builds a \textbf{forward-looking affordability regime
indicator} (FAIR) for UK housing. The maths is deliberately simple ''
it's accounting identities and z-score composites, not econometric
estimation. Here's a structured walkthrough of every mathematical piece,
from the supply identities through to the final FAIR number.

\begin{center}\rule{0.5\linewidth}{0.5pt}\end{center}

\subsection{ðŸ---️ 1. The Supply \& Need
Identities}\label{uxf0uxffuxef-1.-the-supply-need-identities}

These aren't statistical models '' they're \textbf{accounting
frameworks} that express housing as a throughput system where each
factor is a multiplier. If any one term collapses toward zero, the whole
product collapses.

\subsubsection{Identity 1: Affordable Market Supply
(AMS)}\label{identity-1-affordable-market-supply-ams-1}

\[AMS = (HM \cdot P \cdot AR \cdot D \cdot T \cdot PVC) + HS\]

This is a \textbf{multiplicative chain} '' deliverable affordable
housing is the product of:

{\def\LTcaptype{none} % do not increment counter
\begin{longtable}[]{@{}
  >{\raggedright\arraybackslash}p{(\linewidth - 2\tabcolsep) * \real{0.5000}}
  >{\raggedright\arraybackslash}p{(\linewidth - 2\tabcolsep) * \real{0.5000}}@{}}
\toprule\noalign{}
\begin{minipage}[b]{\linewidth}\raggedright
\textbf{Term}
\end{minipage} & \begin{minipage}[b]{\linewidth}\raggedright
\textbf{Meaning}
\end{minipage} \\
\midrule\noalign{}
\endhead
\bottomrule\noalign{}
\endlastfoot
\(HM\) & Housebuilding market capacity (actual deliverable output) \\
\(P\) & Affordable proportion of that delivery \\
\(AR\) & Absorption rate (how fast homes sell without destabilising
price) \\
\(D\) & Diversity of tenure/product mix (widens take-up) \\
\(T\) & Throughput efficiency (permission †' start †' completion) \\
\(PVC\) & Planning value capture that works in practice \\
\(HS\) & Existing affordable housing stock (additive, not
multiplicative) \\
\end{longtable}
}

The key insight: because the first six terms are \textbf{multiplied}, a
bottleneck in \emph{any one} (say absorption rate drops near zero during
a credit crunch) kills the entire flow '' even if land, planning, and
capacity are fine.

\subsubsection{Identity 2: Affordable Housing Need
(AN)}\label{identity-2-affordable-housing-need-an}

Start from the expanded form:

\[AN = HD + (HM \cdot P \cdot AR \cdot D \cdot T \cdot PVC \cdot F) + HS - AMS + NCC\]

Now substitute the AMS definition and simplify. Since
\[AMS = (HM \cdot P \cdot AR \cdot D \cdot T \cdot PVC) + HS\], the
\[HS\] terms cancel and the multiplicative block partially cancels:

\[AN = HD + (HM \cdot P \cdot AR \cdot D \cdot T \cdot PVC \cdot (F - 1)) + NCC\]

\begin{itemize}
\tightlist
\item
  \textbf{\[HD\]} = HomeMaker effective demand (real household need, not
  just ``bankable demand this quarter'')
\item
  \textbf{\[F\]} = sequencing/fast-tracking factor. At \[F = 1\]
  (baseline), the middle term vanishes '' need equals demand plus the
  credit stabiliser. When \[F > 1\], fast-tracking adds supply; when
  \[F < 1\], delays subtract it.
\item
  \textbf{\[NCC\]} = New Circuit of Credit Creation '' a ring-fenced
  funding channel for affordable delivery that doesn't just bid up
  existing stock prices.
\end{itemize}

This is the paper's structural claim in equation form: \emph{need
persists when supply is tethered to private-market throughput, and only
a decoupled credit circuit can break the dependency.}

\begin{center}\rule{0.5\linewidth}{0.5pt}\end{center}

\subsection{ðŸ``Š 2. FAIR '' The Indicator (Step by
Step)}\label{uxf0uxffux161-2.-fair-the-indicator-step-by-step}

FAIR converts the qualitative diagnosis (``markets clear through
rationing, not repricing'') into a \textbf{single reproducible number}
each quarter.

\subsubsection{Step 1: Compute Year-on-Year Growth
Rates}\label{step-1-compute-year-on-year-growth-rates}

For house prices \[P_t\], mortgage stock \[MB_t\], and turnover
\[TO_t\]:

\[g_t^P = \frac{P_t - P\_{t-4}}{P\_{t-4}}\]

\[g_t^{MB} = \frac{MB_t - MB\_{t-4}}{MB\_{t-4}}\]

\[g_t^{TO} = \frac{TO_t - TO\_{t-4}}{TO\_{t-4}}\]

(The subscript \[t-4\] means ``same quarter last year'' in quarterly
data.)

\subsubsection{Step 2: Compute the Credit``Price
Wedge}\label{step-2-compute-the-creditprice-wedge}

\[W_t = g_t^P - g_t^{MB}\]

This is the \textbf{core diagnostic variable}. When \[W_t > 0\], house
prices are growing faster than the mortgage stock that finances them ''
prices are \emph{decoupling} from credit.

\subsubsection{Step 3: Optional New-Build Share
Change}\label{step-3-optional-new-build-share-change}

\[\Delta NB_t = NB_t - NB\_{t-4}\]

A simple year-on-year difference (not a growth rate) in the new-build
share of transactions.

\subsubsection{Step 4: Baseline Normalisation
(Z-Scores)}\label{step-4-baseline-normalisation-z-scores}

Rather than z-scoring against the full history (which would let crisis
periods distort what ``normal'' looks like), it pools two ``relatively
normal'' windows:

\begin{itemize}
\tightlist
\item
  \textbf{2003Q1``2007Q4}
\item
  \textbf{2013Q1``2019Q4}
\end{itemize}

For any series \[X_t\]:

\[z(X_t) = \frac{X_t - \mu_X}{\sigma_X}\]

where \[\mu_X\] and \[\sigma_X\] are computed \textbf{only from the
baseline quarters}.

\subsubsection{Step 5: FAIR Composite}\label{step-5-fair-composite}

\textbf{Three-component version:}

\[FAIR_t = 100 \cdot \Big( 0.55 \cdot z(W_t) - 0.35 \cdot z(g_t^{TO}) + 0.10 \cdot z(\Delta NB_t) \Big)\]

\textbf{Two-component fallback} (when new-build data is unavailable):

\[FAIR_t = 100 \cdot \Big( 0.55 \cdot z(W_t) - 0.35 \cdot z(g_t^{TO}) \Big)\]

\subsubsection{Step 6:
Direction-of-Flow}\label{step-6-direction-of-flow}

\[\Delta FAIR_t = FAIR_t - FAIR\_{t-1}\]

\begin{center}\rule{0.5\linewidth}{0.5pt}\end{center}

\subsection{🎯 3. Interpretation
Bands}\label{uxf0uxffux17e-3.-interpretation-bands}

{\def\LTcaptype{none} % do not increment counter
\begin{longtable}[]{@{}
  >{\raggedright\arraybackslash}p{(\linewidth - 2\tabcolsep) * \real{0.5000}}
  >{\raggedright\arraybackslash}p{(\linewidth - 2\tabcolsep) * \real{0.5000}}@{}}
\toprule\noalign{}
\begin{minipage}[b]{\linewidth}\raggedright
\textbf{FAIR range}
\end{minipage} & \begin{minipage}[b]{\linewidth}\raggedright
\textbf{Regime reading}
\end{minipage} \\
\midrule\noalign{}
\endhead
\bottomrule\noalign{}
\endlastfoot
\(FAIR \ge 50\) & Strong deterioration risk \\
\(20 \le FAIR < 50\) & Mild deterioration bias \\
\(-20 \le FAIR < 20\) & Neutral / noisy \\
\(-50 \le FAIR < -20\) & Mild improvement bias \\
\(FAIR < -50\) & Strong improvement (often itself stress-driven '' e.g.,
post-crash credit collapse) \\
\end{longtable}
}

\begin{center}\rule{0.5\linewidth}{0.5pt}\end{center}

\subsection{ðŸ'' 4. Backtesting
Framework}\label{uxf0uxff-4.-backtesting-framework}

\subsubsection{Crash-Start Definition}\label{crash-start-definition}

A \textbf{crash start} is defined algorithmically as: - A local price
peak followed by a drawdown of at least \[X\\] within \[Y\] quarters -
With a minimum cooldown of \[C\] quarters between events

\subsubsection{Warning Rules
(Illustrative)}\label{warning-rules-illustrative}

{\def\LTcaptype{none} % do not increment counter
\begin{longtable}[]{@{}
  >{\raggedright\arraybackslash}p{(\linewidth - 4\tabcolsep) * \real{0.3333}}
  >{\raggedright\arraybackslash}p{(\linewidth - 4\tabcolsep) * \real{0.3333}}
  >{\raggedright\arraybackslash}p{(\linewidth - 4\tabcolsep) * \real{0.3333}}@{}}
\toprule\noalign{}
\begin{minipage}[b]{\linewidth}\raggedright
\textbf{Rule}
\end{minipage} & \begin{minipage}[b]{\linewidth}\raggedright
\textbf{Condition}
\end{minipage} & \begin{minipage}[b]{\linewidth}\raggedright
\textbf{What it captures}
\end{minipage} \\
\midrule\noalign{}
\endhead
\bottomrule\noalign{}
\endlastfoot
A & \(FAIR > 20\) for 2 consecutive quarters & Sustained elevated
stress \\
B & \(\Delta FAIR > 5\) for 2 consecutive quarters & Sustained
acceleration \\
C & \(FAIR > 0\) and \(\Delta FAIR \ge 0\) & Broad worsening (looser
trigger) \\
\end{longtable}
}

\subsubsection{Evaluation Metrics}\label{evaluation-metrics}

\begin{itemize}
\tightlist
\item
  \textbf{Lead time}
\item
  \textbf{False-positive rate}
\end{itemize}

\begin{center}\rule{0.5\linewidth}{0.5pt}\end{center}

\subsection{ðŸ'¡ 5. Key Takeaways on the
Maths}\label{uxf0uxff-5.-key-takeaways-on-the-maths}

\begin{enumerate}
\def\labelenumi{\arabic{enumi}.}
\tightlist
\item
  \textbf{Supply identities} use multiplication to encode the ``weakest
  link'' property.
\item
  \textbf{FAIR is a weighted z-score composite} with a baseline that
  excludes crisis windows.
\item
  \textbf{No parameters are estimated} '' fixed weights by design.
\item
  \textbf{Baseline choice} is the most consequential methodological
  decision.
\end{enumerate}

\begin{center}\rule{0.5\linewidth}{0.5pt}\end{center}

\protect\phantomsection\label{fair}
\subsection{FAIR Paper}\label{fair-paper}

\begin{center}\rule{0.5\linewidth}{0.5pt}\end{center}

\section{\texorpdfstring{\textbf{Home@ix FAIR:} A Forward Indicator of
UK Housing Affordability Regimes Using Credit``Price Decoupling and
Market
Depth}{Home@ix FAIR: A Forward Indicator of UK Housing Affordability Regimes Using Credit``Price Decoupling and Market Depth}}\label{homeix-fair-a-forward-indicator-of-uk-housing-affordability-regimes-using-creditprice-decoupling-and-market-depth}

\textbf{Author:} Home@ix (Working Paper / Draft)

\textbf{Date:} 25 February 2026

\subsection{Abstract}\label{abstract-1}

UK housing affordability is usually discussed as a price level or a
price-to-income ratio. That framing is necessary but incomplete because
UK housing markets often do not clear primarily through price
adjustment. In stressed regimes, markets clear through quantities and
composition: turnover collapses, chains fail, mortgage-dependent
households are rationed out, and observed prices can appear sticky or
misleading. This paper introduces Home@ix FAIR: a quarterly indicator
designed as a 2``3 year forward regime signal for affordability
understood as access and allocation, not only price burden. FAIR is
intentionally reduced-form and reproducible. It integrates a
credit``price wedge (prices outrunning mortgage-stock growth), a
market-depth term (turnover dynamics) capturing thin-market clearing,
and an optional composition term (new-build share dynamics). To avoid
defining œnormal using crisis regimes and structural breaks, FAIR is
normalised to pooled baseline windows excluding the Global Financial
Crisis aftermath and Covid-era disruptions. FAIR is paired with a
direction-of-flow statistic and can be evaluated via event-based
backtesting (drawdown-defined crash starts). Together, the index and
backtests form an operational toolkit for monitoring affordability
regime risk.

\subsection{Executive Summary}\label{executive-summary-1}

UK housing affordability is typically treated as a static burden ratio.
Yet housing is not a frictionless market. When credit tightens or
uncertainty rises, the UK housing market frequently clears via reduced
participation rather than orderly repricing: transactions fall, turnover
collapses, and access deteriorates even when price ratios appear stable.

This paper introduces \textbf{Home@ix FAIR}, a reproducible quarterly
indicator designed to act as a \textbf{2``3 year forward regime signal}
for affordability understood as \textbf{access and allocation}. FAIR
integrates:

\begin{itemize}
\tightlist
\item
  a \textbf{credit``price wedge} capturing decoupling between
  house-price growth and mortgage-stock growth;
\item
  a \textbf{market-depth / turnover term} capturing thin-market
  (quantity-clearing) stress; and
\item
  an \textbf{optional composition term} based on new-build share
  dynamics.
\end{itemize}

To avoid contaminating baseline moments with crisis regimes and
structural breaks, FAIR is normalised using pooled baseline windows:

\begin{itemize}
\tightlist
\item
  \textbf{Baseline windows:} 2003Q1``2007Q4 and 2013Q1``2019Q4
\item
  \textbf{Excluded:} 2008``2012 (GFC and aftermath) and 2020+ (Covid and
  subsequent regime breaks)
\end{itemize}

FAIR is accompanied by a direction-of-flow measure and can be assessed
as an early-warning tool using objective event definitions (e.g., crash
starts defined by drawdowns from local peaks).

\subsection{Introduction: Affordability as Access, Allocation, and
Market
Clearing}\label{introduction-affordability-as-access-allocation-and-market-clearing}

Standard affordability measures (e.g., price-to-income ratios) treat
affordability as a price-only object. But UK housing often fails to
clear under stress through price adjustment. Instead, the system shifts
toward thin-market clearing:

\begin{itemize}
\tightlist
\item
  transactions decline sharply;
\item
  the transacting population becomes more selected (composition shifts);
\item
  access deteriorates even if headline price ratios appear stable.
\end{itemize}

This paper therefore treats affordability as a \textbf{regime outcome}
shaped by credit creation, market depth, and institutional constraints,
and introduces an indicator designed to detect when the system is moving
into a regime where affordability is likely to deteriorate over the next
2``3 years.

\subsection{From Descriptive Diagnosis to a Forward-Looking
Tool}\label{from-descriptive-diagnosis-to-a-forward-looking-tool}

Home@ix begins from an integrated diagnosis: mortgage-credit dominance
in bank balance sheets, endogenous money creation via bank lending, and
housing™s role as both a necessity and collateral asset. Moving from
diagnosis to a forward-looking tool requires operationalising how
stressed regimes appear in observable data.

\subsubsection{Why œforward-looking is operationally meaningful
here}\label{why-ux153forward-looking-is-operationally-meaningful-here}

A forward-looking affordability signal must capture two patterns that
commonly precede later deterioration:

\begin{itemize}
\tightlist
\item
  \textbf{Credit``price decoupling:} price growth outruns mortgage-stock
  growth.
\item
  \textbf{Thin-market clearing:} market depth collapses (turnover
  declines), indicating rationing and throughput failure.
\end{itemize}

FAIR integrates these mechanisms into a single auditable index with an
explicit baseline definition.

\subsection{Data and Reproducible
Construction}\label{data-and-reproducible-construction}

FAIR is computed from quarterly series produced by the Home@ix feature
set.

\subsubsection{Inputs (quarterly)}\label{inputs-quarterly-1}

Key variables:

\begin{itemize}
\tightlist
\item
  House prices (level): \texttt{avg\_house\_price\_gbp}
\item
  Mortgage stock (level, £m): \texttt{mb\_total\_gbp\_m}
\item
  Turnover (market depth): \texttt{turnover\_pct\_q} (transactions
  scaled by dwelling stock)
\item
  Optional composition proxy: \texttt{newbuild\_share\_of\_transactions}
\item
  Time index: \texttt{period} (quarter), and \texttt{geo} for region
  selection (e.g., EW)
\end{itemize}

\subsubsection{Derived series}\label{derived-series-1}

Define year-on-year (YoY) growth rates:

\[ g_t^P = YoY(P_t) \]

\[ g_t^{MB} = YoY(MB_t) \]

\[ g_t^{TO} = YoY(TO_t) \]

Define the credit``price wedge:

\[ W_t = g_t^P - g_t^{MB} \]

Optional new-build share change (YoY):

\[ \Delta NB_t = NB_t - NB\_{t-4} \]

\subsection{Baseline Normalisation (Core
Proviso)}\label{baseline-normalisation-core-proviso-1}

To avoid defining œnormal using crisis regimes and structural breaks,
z-scores are computed from pooled baseline windows:

\begin{itemize}
\tightlist
\item
  2003Q1``2007Q4
\item
  2013Q1``2019Q4
\end{itemize}

Exclusions:

\begin{itemize}
\tightlist
\item
  2008``2012 (Global Financial Crisis and aftermath)
\item
  2020+ (Covid and subsequent regime breaks)
\end{itemize}

For any series \textbackslash X\_t\textbackslash, define:

\[ z(X_t) = \frac{X_t - \mu_X}{\sigma_X} \]

where \[\mu_X\] and \[\sigma_X\] are computed using baseline quarters
only.

\subsection{The Home@ix FAIR Indicator}\label{the-homeix-fair-indicator}

\subsubsection{FAIR definition}\label{fair-definition-1}

Default (three-component) specification:

\[ FAIR_t = 100 \cdot \Big( 0.55 \cdot z(W_t) - 0.35 \cdot z(g_t^{TO}) + 0.10 \cdot z(\Delta NB_t) \Big) \]

If the new-build term is not available, the final term is omitted:

\[ FAIR_t = 100 \cdot \Big( 0.55 \cdot z(W_t) - 0.35 \cdot z(g_t^{TO}) \Big) \]

\subsubsection{Direction-of-flow}\label{direction-of-flow-1}

Define the direction-of-flow statistic:

\[ \Delta FAIR_t = FAIR_t - FAIR\_{t-1} \]

This is used for monitoring whether affordability pressure is
accelerating or easing.

\subsubsection{Interpretation bands
(practical)}\label{interpretation-bands-practical}

A simple interpretive scheme:

\begin{itemize}
\tightlist
\item
  \[FAIR \ge 50\]: strong deterioration regime risk
\item
  \[20 \le FAIR < 50\]: mild deterioration bias
\item
  \[-20 \le FAIR < 20\]: neutral/noisy
\item
  \[-50 \le FAIR < -20\]: mild improvement bias
\item
  \[FAIR < -50\]: strong improvement regime (often stress-driven)
\end{itemize}

\subsection{Using FAIR for Forward-Looking Evaluation (Event
Backtesting)}\label{using-fair-for-forward-looking-evaluation-event-backtesting}

FAIR can be evaluated using objective events and simple warning rules.

\subsubsection{Crash-start events
(rule-based)}\label{crash-start-events-rule-based-1}

A œcrash start may be defined as a local price peak followed by a
drawdown of at least \[X\\] within \[Y\] quarters, with a minimum
separation (cooldown) of \[C\] quarters to avoid double-counting
clustered peaks.

\subsubsection{Illustrative warning
rules}\label{illustrative-warning-rules-1}

Examples:

\begin{itemize}
\tightlist
\item
  A: sustained level stress: \[FAIR > 20\] for 2 quarters
\item
  B: sustained acceleration: \[\Delta FAIR > 5\] for 2 quarters
\item
  C: broad stress-worsening: \[FAIR > 0\] and \[\Delta FAIR \ge 0\]
\end{itemize}

\subsubsection{Metrics}\label{metrics-1}

\begin{itemize}
\tightlist
\item
  Lead time (quarters) from the most recent signal to each crash start.
\item
  False-positive proxy: share of signal quarters not followed by any
  crash start within a forward window (e.g., 8 quarters).
\end{itemize}

\subsection{Outputs and Visualisations
(Reproducible)}\label{outputs-and-visualisations-reproducible-1}

The FAIR build produces:

\begin{itemize}
\tightlist
\item
  a FAIR level figure (headline chart),
\item
  a component contributions figure (mechanism chart),
\item
  an animated direction-of-flow GIF emphasising \[\Delta FAIR_t\],
\item
  and a tidy CSV export of intermediate series for auditability.
\end{itemize}

\subsubsection{Figure placeholders}\label{figure-placeholders}

Replace the filenames with the paths to the generated assets.

Figure 1. Home@ix FAIR (EW): Level. Baseline windows: 2003Q1``2007Q4 and
2013Q1``2019Q4; excludes 2008``2012 and 2020+.

Figure 2. Home@ix FAIR (EW): Component contributions. FAIR is the sum of
wedge, turnover, and optional new-build terms.

\subsection{Limitations and
Extensions}\label{limitations-and-extensions-1}

FAIR is a reduced-form indicator derived from aggregate market series.
It does not directly encode distributional constraints (income
heterogeneity, cash-buyer share, arrears microdata) or direct developer
pacing series (e.g., sales per outlet per week) unless explicitly added.
These are natural extensions; the baseline FAIR specification is
intentionally reproducible and captures two dominant UK mechanisms:
credit``price decoupling and thin-market clearing.

\subsection{Conclusion}\label{conclusion-1}

UK housing affordability is not only a price level problem; it is also
an allocation and throughput problem. Home@ix FAIR provides a
transparent, reproducible way to detect when the system is shifting into
a regime where affordability is likely to deteriorate over a 2``3 year
horizon via credit``price decoupling and/or collapsing market depth.
Combined with direction-of-flow and event backtesting, FAIR is
positioned as an operational forward-looking tool: a single œNumber
grounded in auditable data and interpretable mechanisms.

\subsection{Reproducibility Note}\label{reproducibility-note-1}

The figures referenced in Section 8 are produced by the FAIR build
scripts in the Home@ix workflow. The output directory should contain:

\begin{itemize}
\tightlist
\item
  \texttt{fair\_quarterly\_audit.csv}
\item
  \texttt{fig\_fair\_level.png}
\item
  \texttt{fig\_fair\_contributions.png}
\item
  \texttt{anim\_fair\_direction\_of\_flow.gif}
\end{itemize}

\begin{center}\rule{0.5\linewidth}{0.5pt}\end{center}

\begin{center}\rule{0.5\linewidth}{0.5pt}\end{center}

\end{document}
